% Options for packages loaded elsewhere
\PassOptionsToPackage{unicode}{hyperref}
\PassOptionsToPackage{hyphens}{url}
\documentclass[
]{article}
\usepackage{xcolor}
\usepackage{amsmath,amssymb}
\setcounter{secnumdepth}{-\maxdimen} % remove section numbering
\usepackage{iftex}
\ifPDFTeX
  \usepackage[T1]{fontenc}
  \usepackage[utf8]{inputenc}
  \usepackage{textcomp} % provide euro and other symbols
\else % if luatex or xetex
  \usepackage{unicode-math} % this also loads fontspec
  \defaultfontfeatures{Scale=MatchLowercase}
  \defaultfontfeatures[\rmfamily]{Ligatures=TeX,Scale=1}
\fi
\usepackage{lmodern}
\ifPDFTeX\else
  % xetex/luatex font selection
\fi
% Use upquote if available, for straight quotes in verbatim environments
\IfFileExists{upquote.sty}{\usepackage{upquote}}{}
\IfFileExists{microtype.sty}{% use microtype if available
  \usepackage[]{microtype}
  \UseMicrotypeSet[protrusion]{basicmath} % disable protrusion for tt fonts
}{}
\makeatletter
\@ifundefined{KOMAClassName}{% if non-KOMA class
  \IfFileExists{parskip.sty}{%
    \usepackage{parskip}
  }{% else
    \setlength{\parindent}{0pt}
    \setlength{\parskip}{6pt plus 2pt minus 1pt}}
}{% if KOMA class
  \KOMAoptions{parskip=half}}
\makeatother
\usepackage{longtable,booktabs,array}
\usepackage{calc} % for calculating minipage widths
% Correct order of tables after \paragraph or \subparagraph
\usepackage{etoolbox}
\makeatletter
\patchcmd\longtable{\par}{\if@noskipsec\mbox{}\fi\par}{}{}
\makeatother
% Allow footnotes in longtable head/foot
\IfFileExists{footnotehyper.sty}{\usepackage{footnotehyper}}{\usepackage{footnote}}
\makesavenoteenv{longtable}
\setlength{\emergencystretch}{3em} % prevent overfull lines
\providecommand{\tightlist}{%
  \setlength{\itemsep}{0pt}\setlength{\parskip}{0pt}}
\usepackage{newunicodechar}
\newunicodechar{→}{\ensuremath{\rightarrow}}
\newunicodechar{≈}{\ensuremath{\approx}}
\newunicodechar{≤}{\ensuremath{\leq}}
\newunicodechar{−}{\ensuremath{-}}
\newunicodechar{×}{\ensuremath{\times}}
\newunicodechar{Δ}{\ensuremath{\Delta}}
\newunicodechar{µ}{\textmu}
\usepackage{bookmark}
\IfFileExists{xurl.sty}{\usepackage{xurl}}{} % add URL line breaks if available
\urlstyle{same}
\hypersetup{
  hidelinks,
  pdfcreator={LaTeX via pandoc}}

\author{}
\date{}

\begin{document}

\section{quadbit: Hand-Written FP4 Tensor-Core Kernels and a Deployment
Stack for Consumer Blackwell
(SM120)}\label{quadbit-hand-written-fp4-tensor-core-kernels-and-a-deployment-stack-for-consumer-blackwell-sm120}

\textbf{Draft v0.1.} Working paper. Every quantitative claim is measured
on a Modal cloud RTX PRO 6000 (SM120, no tcgen05) unless tagged
otherwise, and traces to a harness in \texttt{harness/} or a kernel in
\texttt{cuda/}. Numbers here are copied from
\texttt{docs/paper\_notes.md} and \texttt{docs/standing.md}; if they
disagree, those two files are the source of truth and this draft is
stale.

\begin{center}\rule{0.5\linewidth}{0.5pt}\end{center}

\subsection{Abstract}\label{abstract}

Consumer and pro Blackwell cards (SM120: RTX PRO 6000, RTX 5090) ship
FP4 tensor cores, and the dense-NVFP4 software that reaches them is now
strong. FlashInfer's \texttt{mm\_fp4} (\texttt{b12x}/\texttt{cutlass}),
vLLM's \texttt{modelopt\_fp4}, and SGLang's FlashInfer CUTLASS
\texttt{fp4\_gemm} all bind the native W4A4 path on this card, and on a
same-card, same-shape, fp32-reference-gated leaderboard they beat our
hand-written two-level dense kernel by 1.35 to 2.2x. We report that
plainly: quadbit does not claim dense FP4 speed leadership on SM120.
Dense FP4 is commoditized; the open Pareto corner is \emph{sparse} FP4,
and that is what quadbit targets. (SM120 also lacks the
\texttt{tcgen05}/UMMA tensor-memory path the datacenter B200 (SM100)
uses, so every FP4 GEMM here runs through warp-level \texttt{mma.sync}
and \texttt{mma.sp}, whose operand, scale, and 2:4-metadata bit-layouts
we first derived by probe-and-verify, validated to relative error 0.)

We contribute a deployable 2:4-sparse FP4 stack for SM120.
\textbf{First, two-level sparse FP4 kernels} (per-16 \texttt{ue4m3}
local plus per-row and per-column fp32 global rescale) whose deployed
accuracy equals the trained fake-quant/QAT checkpoint, closing the
single-level deploy gap (11.89 → 8.95 PPL); they beat CUTLASS 80b (the
only other sparse FP4 kernel in existence) on every shape and beat even
the best dense FlashInfer kernel in wall-clock on every Llama-3-8B
prefill shape (1.07 to 1.38x) --- if a weight can be 2:4-pruned, quadbit
sparse is the fastest way to run its FP4 GEMM on the platform.
\textbf{Second, serving integration}: the sparse MLP runs inside a
vLLM-style stack as a \texttt{torch.library} custom op that CUDA-graph
capture includes, giving a real end-to-end Pareto row on recovered
Llama-3.1-8B-Instruct --- sparse wins decode and total request latency
in the interactive/low-batch, long-generation regime (81 of 112 cells)
while NVFP4 keeps the prefill-bound corner. \textbf{Third,
cross-architecture sparse-policy transfer}: on DeepSeek-V4-Flash and
GLM-5.2-NVFP4 (whose Deep Sparse Attention runs natively on SM120) the
accuracy cost of sparsity is a \emph{placement} problem ---
down-projection sparsity is nearly free while gate/up carries the tax,
and a route-slot policy (top experts dense, low-weight tail 2:4) is the
best quality/sparse-FLOP tradeoff. On GLM-5.2 the DeepSeek rule
transfers: route-slot D2 costs +0.065 held-out PPL and preserves a
4-task downstream smoke-suite average to within about one point of dense
(.7508 vs .7603, no task collapsing). The same down-anchor rule
transfers to a third architecture, \textbf{gpt-oss-20b}, which we serve
end-to-end in vLLM 0.24 with training-free recovery (all-expert 2:4 is
incoherent at 111 PPL; anchoring gate/up dense recovers to 6.499, within
1.66x of stock) and a fused monolithic sparse MoE (+6\% serve). On
accuracy the deployed dense path is W4A4 and costs +0.63 PPL with no
calibration; sparse deploys at its trained accuracy but stays a real PPL
behind dense, so sparse is a speed Pareto point conditioned on
prunability, not an accuracy win.

The deployed sparse MoE policy path is graph-enabled on SM120 (P4), and
the prior dense anchored/grouped projection bottleneck is removed by
delegating those projections to FlashInfer's native grouped NVFP4 GEMM
(\texttt{group\_gemm\_nvfp4\_nt\_groupwise}, C1), with no custom dense
grouped-GEMM required. DeepSeek-D2 and GLM route-slot D2 now run as
graph-enabled, quality-matching deployed policies with native SM120 DSA;
native-delegate captured DeepSeek-D2 decodes faster than eager (5.82 vs
4.04 tok/s at matched PPL). Limitations we state up front: a direct
same-harness board (C2) shows the dense NVFP4 fused MoE decodes 6-8x
faster than the sparse policy and uses less memory, so this is a
quality-preservation and graph-transfer result, not a MoE decode-speed
SOTA; the C1 speedup is a same-model/same-policy Pareto result against
our own dequant-loop and eager paths, not a production-wide decode-speed
win, and dense FP4 speed belongs to the ecosystem baselines; the native
delegate depends on FlashInfer availability and its swizzled NVFP4 scale
layout; the GLM graph run is validated on a short held-out passage;
GLM-5.2 requires 8x RTX PRO 6000; the GLM downstream evidence is a small
4-task smoke suite, not an exhaustive benchmark; and all-MLP sparsity
carries a real PPL tax that neither training-free repair nor any
layerwise QAT objective closes (three QAT objectives, from per-expert to
global router-weighted MoE QAT, run to the ground -- all negative, the
tightest reconstruction giving the worst downstream).

\begin{center}\rule{0.5\linewidth}{0.5pt}\end{center}

\subsection{1. Introduction}\label{introduction}

FP4 is the smallest numeric format with tensor-core support on
Blackwell, and on paper it promises roughly 4x the throughput of bf16
and 4x smaller weights. When we began, the SM120 FP4 software was thin
and the natural target was a fast dense kernel. That target has since
been commoditized: FlashInfer, vLLM, and SGLang now all reach the native
dense NVFP4 W4A4 path on this card, and FlashInfer's CUDA-13
\texttt{b12x}/\texttt{cutlass} kernels beat our hand-written dense
kernel by 1.35 to 2.2x (Section 4). Dense FP4 on SM120 is no longer
where the open problem is.

The open Pareto corner is \emph{sparse} FP4. No mainstream serving stack
exposes a 2:4-sparse FP4 deployment path on SM120: FlashInfer, SGLang,
and vLLM ship none, and CUTLASS 80b --- the only other sparse FP4 kernel
in existence --- is an unwrapped example with documented block-scaled
problems, not a model-level deployment. quadbit's differentiator is
therefore not dense speed but \emph{deployable} sparse FP4: a kernel
whose semantics match fake-quant/QAT, wired into a real serving stack,
with model-level sparsity policy validated across Llama, DeepSeek, and
GLM-style architectures. This paper measures that corner honestly ---
where sparse FP4 wins, where it does not, and what the accuracy costs
once the hardware quantizes activations too. Answering it required
deriving the SM120 FP4 operand, scale, and 2:4-metadata layouts
ourselves (undocumented at the bit level), building the sparse kernel to
the bandwidth roofline, and quantifying accuracy on real checkpoints
rather than asserting it.

Our contributions:

\textbf{C1. An SM120 FP4 landscape measurement (Sections 4-5).} On a
same-card, same-shape, fp32-reference-gated backend leaderboard, quadbit
dense loses to FlashInfer best-of-backend
(\texttt{b12x}/\texttt{cutlass}) by 1.35 to 2.2x, so we pivot away from
any dense-speed-leadership claim. The sparse kernel, by contrast, beats
CUTLASS 80b (the only other sparse FP4 kernel) on every shape and beats
even the best dense FlashInfer kernel in wall-clock on every Llama-3-8B
prefill shape (1.07 to 1.38x) --- the SM120 sparse-FP4 operand, scale,
and 2:4-metadata layouts derived by probe-and-verify (relative error 0)
underpin it.

\textbf{C2. Deployable two-level sparse FP4 (Sections 3, 5, 8).} A
two-level sparse kernel (per-16 \texttt{ue4m3} local plus per-row and
per-column fp32 global rescale) whose deployed accuracy \emph{equals}
the trained fake-quant/QAT checkpoint, closing the single-level deploy
gap (11.89 → 8.95 PPL) that made prior sparse-FP4 results undeployable.
Around it: a packer, fused activation quantizer, \texttt{nn.Linear}
drop-in, and fused block kernels.

\textbf{C3. Serving integration (Section 9).} The sparse MLP runs inside
a vLLM-style serving stack as a \texttt{torch.library} custom op that
vLLM's fullgraph compile and CUDA-graph capture include, with correct
sparse output (guarded against a silent dense fall-back). On recovered
Llama-3.1-8B-Instruct, graph-vs-graph, sparse wins decode and total
request latency in the interactive/low-batch, long-generation regime (81
of 112 cells); NVFP4 keeps the prefill-bound corner. The MoE plugin path
(C4) now graph-captures for the deployed sparse policies (P4), and C1
removes the dense-anchor decode bottleneck by delegating the
anchored/grouped projections to FlashInfer's native grouped NVFP4 GEMM
(\texttt{group\_gemm\_nvfp4\_nt\_groupwise}), with no custom dense
grouped-GEMM, so the captured DeepSeek-D2 path now decodes faster than
eager at matched quality.

\textbf{C4. Cross-architecture sparse-policy transfer (Section 10).} On
DeepSeek-V4-Flash and GLM-5.2 the accuracy cost of sparsity is a
\emph{placement} problem that transfers across architectures:
down-projection sparsity is far less damaging than gate/up, and a
route-slot policy (top experts dense, low-weight tail 2:4) gives the
best quality/sparse-FLOP tradeoff. A 4-task GLM downstream smoke suite
rebuts the ``PPL says fine but downstream collapses'' concern for the
route-slot D2 policy (AVG .7508 vs dense .7603, no task collapsing). The
rule reaches a \textbf{third architecture}, gpt-oss-20b, served
end-to-end in vLLM with training-free recovery (111 -\textgreater{}
6.499 PPL) and a fused monolithic sparse MoE (Section 10.3); and we
close the ``can training repair the gate/up tax?'' question in the
negative across the whole layerwise family (per-expert and global
router-weighted MoE QAT both fail, tightest reconstruction giving worst
downstream, Section 10.2).

\begin{center}\rule{0.5\linewidth}{0.5pt}\end{center}

\subsection{2. Background: FP4 on SM120}\label{background-fp4-on-sm120}

\textbf{The formats.} NVFP4 encodes values as E2M1 (the 4-bit
\texttt{\{0,\ .5,\ 1,\ 1.5,\ 2,\ 3,\ 4,\ 6\}} grid with sign) in blocks
of 16, with a two-level scale: a per-16 local scale in \texttt{ue4m3}
and a per-row fp32 global scale, where the global rescales the local
codes into \texttt{ue4m3}'s precise range (global = rowamax / 2688, and
2688 = 448 x 6 is the product of the \texttt{ue4m3} and E2M1 maxima).
MXFP4 uses blocks of 32 with a single power-of-two \texttt{ue8m0} scale;
because that scale carries no mantissa, it has no per-block rounding
bias, which matters for accuracy (Section 7).

\textbf{What the tensor core runs.} The FP4 mma multiplies fp4 by fp4.
There is no weight-only FP4 GEMM in silicon. So the deployed kernel is
always W4A4: both operands are 4-bit. Any accuracy number that keeps
activations at 16-bit (W4A16) describes a kernel that does not exist on
this hardware; it is a dequant-then-bf16-matmul path (Marlin-class), not
an FP4 GEMM. This distinction drives the entire accuracy discussion.

\textbf{No tcgen05.} SM100 (B200) accumulates in tensor memory via UMMA.
SM120 has none of that; the only route to the FP4 tensor cores is
warp-level \texttt{mma.sync} (dense) and \texttt{mma.sp} (2:4-sparse),
with accumulators living in the register file. This is the structural
constraint behind every kernel decision below.

\textbf{Measured hardware ceilings (RTX PRO 6000).} Register-only mma
peaks (dead-code-elimination defeated): sparse \texttt{mma.sp} m16n8k128
= 3626k GFLOP/s, dense \texttt{mma.sync} m16n8k64 = 1811k, a ratio of
2.002x (2:4 is a real 2x FLOP feature). L2-to-smem TMA bandwidth ceiling
= 7.3 TB/s. Peak DRAM bandwidth = 1.46 TB/s. Shared memory = 99 KB per
block, 100 KB per SM; L2 = 128 MB. Dense FP4 is compute-bound at roughly
84\% of the mma peak; sparse FP4 is load-bound against a swizzle floor
near 6.0 TB/s.

\begin{center}\rule{0.5\linewidth}{0.5pt}\end{center}

\subsection{3. Deriving the FP4 layouts}\label{deriving-the-fp4-layouts}

None of the SM120 FP4 bit-layouts we needed are documented, so we
recovered them by constructing inputs with known structure, running the
mma, and checking which arrangement reproduced the reference to relative
error 0.

\begin{itemize}
\tightlist
\item
  \textbf{Scale layout.} For the dense block-scaled mma,
  \texttt{scaleA{[}row\ r{]}{[}kb{]}} maps to lane
  \texttt{(r\ \&\ 7)\ *\ 4\ +\ (r\ \textgreater{}\textgreater{}\ 3)},
  byte \texttt{kb}; \texttt{scaleB{[}col\ c{]}{[}kb{]}} maps to lane
  \texttt{c\ *\ 4}, byte \texttt{kb}.
\item
  \textbf{2:4 metadata.} For \texttt{mma.sp}, lane \texttt{L} maps to
  mma-row
  \texttt{(L\ \&\ 1)\ *\ 8\ +\ (L\ \textgreater{}\textgreater{}\ 2)},
  half \texttt{(L\ \textgreater{}\textgreater{}\ 1)\ \&\ 1}, with the
  kept-pair nibble encoded
  \texttt{idx0\ \textbar{}\ (idx1\ \textless{}\textless{}\ 2)}.
\item
  \textbf{NVFP4 dense scale.} The \texttt{scale\_vec::4X} \texttt{ue4m3}
  case uses the same A/B row-to-lane mapping as the \texttt{ue8m0}
  \texttt{2X} case, just with a 4-byte (four per-16) scale register. We
  confirmed this with a wide-range scale probe (2\^{}-4 to 2\^{}2),
  which is the strong test: a uniform-scale probe cannot tell whether
  the four scale bytes map to the correct k-sub-blocks; the wide-range
  one can. Result: relative error 0.
\end{itemize}

The \texttt{ldmatrix} story is a negative result worth recording. The
only sub-byte \texttt{ldmatrix} ptxas accepts on \texttt{sm\_120a} is
the format-converting \texttt{m8n16}/\texttt{m16n16\ .b8x16.b4x16\_p64},
which expands each 4-bit value into its own byte. Our mma wants packed
\texttt{e2m1x2}, so the expanded fragment is the wrong shape and would
need 2x the registers to hold a full operand, which is fatal at the
255-register ceiling. We patched CubeCL to emit the instruction and
confirmed it assembles and runs; it is simply the wrong tool for this
kernel.

\begin{center}\rule{0.5\linewidth}{0.5pt}\end{center}

\subsection{4. Dense FP4 to the compute
ceiling}\label{dense-fp4-to-the-compute-ceiling}

The dense kernel stages packed \texttt{e2m1x2} tiles through shared
memory and issues a wide 2x8 grid of \texttt{mma.sync} instructions per
warp, feeding 16 independent f32 accumulator chains back in as C so
products accumulate across K. Those 16 chains are the instruction-level
parallelism that hides the FP4 mma (OMMA) latency; the kernel is
latency-and-ILP-bound, not occupancy-bound. Widening the warp tile from
2x2 to 2x4 to 2x8 climbed roughly 160k, 230k, 253k GFLOP/s, and 2x8 uses
exactly 255 registers per thread with zero spills, saturating the
register file (8 warps x 255 = 65280 of 65536 registers per SM). This is
the classic ``better performance at lower occupancy'' regime.

\textbf{The false-roofline lesson.} An early mem-only probe suggested a
ceiling that turned out to be an artifact of the probe's own narrow load
pattern, which extracted only 54\% of the L2-to-smem bandwidth. The
lesson, which recurs throughout the project: never declare a memory
roofline from a probe whose own load pattern is the limit.

\textbf{Versus CUTLASS 79b.} On square sizes, apples-to-apples (both
cudaEvent-timed over 20 iterations, no torch dispatch), our dense kernel
reaches 758 / 1220 / 1510 TF/s at 2048 / 4096 / 8192 versus CUTLASS
79b's 634 / 1222 / 1497: a win, a tie, and a win. But the square win
does not generalize. On the rectangular Llama-3-8B GEMM shapes that
actually run, we lose to 79b: attention 4096-cubed at 0.89x, ffn-up
(N=14336) at 0.93x, ffn-down (K=14336) at 1.01x, all 79b-verified.
CUTLASS's tile and schedule adapt to rectangular shapes better than our
fixed tiling. Against 79b specifically the dense kernel is
``competitive, slightly behind on shipping shapes'' while delivering 3.0
to 3.7x over cuBLAS bf16, but 79b is no longer the baseline that
matters: the leaderboard below shows FlashInfer beating us outright, so
dense is not where we win.

We built and measured three persistent/stream-K variants to try to beat
data-parallel on the rectangular shapes (stream-K, a true CUTLASS-style
persistent cross-tile pipeline, and split-K via an f32 global
reduction). All three regressed. The root cause is register pressure:
the 128-register FP4 accumulator tile plus a software scheduler's cursor
state pins the kernel near 254 registers and throttles the mma stream,
while the hardware block scheduler already overlaps consecutive tiles
for free. On SM120 FP4, the huge accumulator tile leaves no register
headroom for a software scheduler, so data-parallel is at the practical
ceiling.

\textbf{The deployed W4A4 kernel is block-scaled, and its scale loads
were an exposed stall.} The 758/1220/1510 figures are the unit-scale
kernel (compile-time-zero scales). The kernel that actually ships for
accuracy is the two-level NVFP4 variant (Section 7), which must stream
per-16 \texttt{ue4m3} scales from global memory each K-step; on SM120
those scales are 8-byte-pitch, so they cannot ride the tile TMA (which
requires 16-byte strides) and were loaded synchronously between the tile
\texttt{try\_wait} and the mma, exposing \textasciitilde500 cycles of
latency every step. Double-buffering the scales and prefetching one step
ahead with \texttt{cp.async} (while STAGES=2 and the 128-wide tile stay
fixed) overlaps that load with the previous step's mma and lifts the
deployed kernel 1.08 to 1.22x (square 8192 865 to 1055 TF/s, biggest
where the k-step count is largest), at relative error 0 against the
synchronous version. This narrows the block-scaled-vs-unit-scale gap
without changing the mma; it is orthogonal to the 758/1220/1510
unit-scale ceiling above.

\textbf{The full SM120 FP4 backend leaderboard (and why we lose the
dense race).} CUTLASS 79b is no longer the only competitor. FlashInfer's
\texttt{mm\_fp4} now exposes several NVFP4 GEMM backends, and on SM120
its \texttt{auto} prefers \texttt{b12x} (a CUDA-13-only SM120/121
kernel), then \texttt{cutlass}, then \texttt{cudnn} (\texttt{trtllm} and
\texttt{cute-dsl} refuse SM120 outright). We benchmarked the deployed
two-level quadbit kernel against every FlashInfer backend on one RTX PRO
6000, over a shape suite spanning square (2048 to 16384), Llama-3-8B
prefill, decode small-M (1 to 128), and serving M = B\emph{S, with an
identical correctness gate: each backend quantizes the same bf16
operands with its own native quantizer and its output is scored against
the fp32 reference (cosine \textgreater{} 0.97 to count; all
NVFP4-on-random-Gaussian rows land at cos 0.991, and the }identical*
cos/maxrel/mae across backends cross-validates that they compute the
same math). Best-backend-per-shape, GEMM-only (effective TF/s =
2\emph{M}N*K / wall):

\begin{longtable}[]{@{}llll@{}}
\toprule\noalign{}
shape & quadbit dense & FlashInfer best (backend) & FI / quadbit \\
\midrule\noalign{}
\endhead
\bottomrule\noalign{}
\endlastfoot
square 8192 & 1045 & 1433 (cutlass) & 1.37x \\
prefill attn 4096-cubed & 838 & 1283 (b12x) & 1.53x \\
prefill ffn-up (N=14336) & 936 & 1374 (auto) & 1.47x \\
prefill ffn-down (K=14336) & 1017 & 1408 (b12x) & 1.38x \\
serving ffn-down (M=65536) & 639 & 1416 (cutlass) & 2.22x \\
\end{longtable}

FlashInfer's \texttt{b12x}/\texttt{cutlass} NVFP4 kernels beat our
hand-written two-level dense by 1.35 to 2.2x. The honest dense story is
no longer ``competitive with CUTLASS 79b''; the SM120 dense-FP4 baseline
moved up and quadbit does not win it. Two structural findings from the
same run: (1) FlashInfer's \texttt{cudnn} backend returns
\texttt{No\ execution\ plans\ support\ the\ graph} on \emph{every}
shape, because the shipped cuDNN (9.10) is below the 9.14 that SM120 FP4
needs; (2) \texttt{b12x} collapses from \textasciitilde1400 to
\textasciitilde640 TF/s once M \textgreater= 65536 (large serving
batch), where only the \texttt{cutlass} backend holds
\textasciitilde1400, so on SM120 the fast dense backend is M-dependent.
One toolchain fact frames the whole comparison: \texttt{b12x} requires
CUDA 13, while quadbit's \texttt{sm\_120a} block-scale mma
(\texttt{kind::mxf4nvf4}/\texttt{block\_scale}/\texttt{scale\_vec::4X})
is \emph{rejected} by ptxas 13 and only assembles under CUDA \textless=
12.8, so the two kernels cannot coexist in one container. Where quadbit
wins is not dense; it is sparse (Section 5).

\begin{center}\rule{0.5\linewidth}{0.5pt}\end{center}

\subsection{5. Sparse FP4 to the bandwidth
roofline}\label{sparse-fp4-to-the-bandwidth-roofline}

Sparse FP4 is load-bound, so the kernel win is a memory-traffic win. The
design stages arbitrary per-group 2:4 metadata and real per-block
\texttt{ue4m3} scales coalesced through a full/empty async pipeline (no
CTA-wide \texttt{\_\_syncthreads}), on a shared-B 256x128
traffic-optimal tiling.

\textbf{The wide-TMA-plus-swizzle breakthrough.} A mem-only probe
suggested a ``2012k roofline''; it was false for the same reason as
Section 4, narrow TMA boxes extracted only 54\% of the 7.3 TB/s
L2-to-smem ceiling. Loading WK=2 k128-slices per TMA (wider boxes) hits
the ceiling, but wide smem rows cause \texttt{ldmatrix} bank conflicts.
Swizzled TMA fixes that (A box 64B with \texttt{SWIZZLE\_64B} and an
\texttt{ldmatrix} XOR
\texttt{off\ \^{}=\ ((off\ \textgreater{}\textgreater{}\ 7)\ \&\ 3)\ \textless{}\textless{}\ 4};
B box 128B with \texttt{SWIZZLE\_128B} and XOR
\texttt{((off\ \textgreater{}\textgreater{}\ 7)\ \&\ 7)\ \textless{}\textless{}\ 4}).
Result: unit-scale sparse went from 2012k to 2731k (+36\%), and the
deployable kernel from 1486k to 2116k (+42\%).

\textbf{Versus CUTLASS 80b, the only other sparse FP4 kernel.} No
shipping library provides a sparse FP4 GEMM on SM120: FlashInfer,
SGLang, and vLLM expose dense NVFP4 and MXFP4-MoE paths only, and
CUTLASS's 80b is an unwrapped example. quadbit is therefore the only
\emph{deployed} 2:4-sparse FP4 GEMM on the platform. Against 80b,
correctness-gated on 80b's own reference check (which passes at every
size, so CUTLASS issue \#3096's block-scaled bug does not affect this
comparison), the deployed two-level sparse kernel wins on every shape:
square 8192 at 1.09x (1973 vs 1807 TF/s), attention 4096-cubed at 1.08x,
ffn-up at 1.01x, ffn-down at 1.12x, all 80b-verified.

\textbf{The Pareto point: sparse beats the best available \emph{dense}
FP4.} The reviewer-obvious result is cross-table. On the same card and
shapes, quadbit's two-level sparse kernel beats not just its own dense
and CUTLASS 80b, but the \emph{fastest FlashInfer dense backend} in
wall-clock on every Llama-3-8B prefill shape: attention 4096-cubed 0.100
vs 0.107 ms (1.07x), ffn-up 0.301 vs 0.350 ms (1.16x), ffn-down 0.265 vs
0.342 ms (1.29x), square 8192 0.557 vs 0.767 ms (1.38x). So if a weight
can be 2:4-pruned, the fastest way to run its FP4 GEMM on SM120 is
quadbit sparse, faster than the best dense FP4 kernel anyone ships. That
is the Pareto point no library currently provides. (Effective FLOP
counts the pruned zeros; the honest comparator is wall-clock for a fixed
GEMM problem, which is what these ratios are. The two sides were
measured in separate containers, CUDA 12.8 for sparse and CUDA 13 for
FlashInfer, forced by the mma/ptxas split noted in Section 4, on the
same physical card at the same clocks.)

\textbf{Ceiling honesty.} Sparse's real advantage over our own dense FP4
is roughly 1.33x at the roofline (2012k vs 1510k deployable at 8192),
not the 2x the mma FLOP ratio suggests. The hardware 2x is a
datacenter-bandwidth feature we cannot reach on SM120. At 1.33x, the
accuracy cost of sparsity has to be small to justify the recovery
pipeline over dense, which is exactly the tension Section 8 confronts.

\textbf{Throughput summary (M=N=K, vs cuBLAS bf16), speed-path ceiling:}

\begin{longtable}[]{@{}
  >{\raggedright\arraybackslash}p{(\linewidth - 8\tabcolsep) * \real{0.0811}}
  >{\raggedright\arraybackslash}p{(\linewidth - 8\tabcolsep) * \real{0.1757}}
  >{\raggedright\arraybackslash}p{(\linewidth - 8\tabcolsep) * \real{0.1757}}
  >{\raggedright\arraybackslash}p{(\linewidth - 8\tabcolsep) * \real{0.2432}}
  >{\raggedright\arraybackslash}p{(\linewidth - 8\tabcolsep) * \real{0.3243}}@{}}
\toprule\noalign{}
\begin{minipage}[b]{\linewidth}\raggedright
size
\end{minipage} & \begin{minipage}[b]{\linewidth}\raggedright
cuBLAS bf16
\end{minipage} & \begin{minipage}[b]{\linewidth}\raggedright
CUTLASS FP4
\end{minipage} & \begin{minipage}[b]{\linewidth}\raggedright
dense FP4 (ours)
\end{minipage} & \begin{minipage}[b]{\linewidth}\raggedright
2:4-sparse FP4 (ours)
\end{minipage} \\
\midrule\noalign{}
\endhead
\bottomrule\noalign{}
\endlastfoot
4096 & 372 TF/s & 1222 & 1136 (3.06x bf16) & 1512 (4.07x bf16) \\
8192 & 423 TF/s & 1497 & 1556 (3.68x bf16) & 2207 (5.22x bf16) \\
16384 & 405 TF/s & n/a & 1645 (4.06x bf16) & 1782 (4.39x bf16) \\
\end{longtable}

These are the \textbf{speed-path ceiling} numbers: MXFP4-fast dense
(\texttt{ue8m0} per-32) and unit-scale sparse, the fastest variants, not
the accuracy-deployed two-level kernels. The \textbf{deployed two-level}
kernels that carry the accuracy recipe are slower and are the ones on
the leaderboard in Section 4 (dense square-8192 1045, not 1556) and the
sparse head-to-head above (sparse square-8192 1973, not 2207). Reporting
both protocols matters: the two-level rescale that buys the accuracy
result costs throughput, so the deployed sparse win over the best
FlashInfer dense (1.07 to 1.38x) is measured on the \emph{deployed}
kernel, not this ceiling.

\begin{center}\rule{0.5\linewidth}{0.5pt}\end{center}

\subsection{6. Deployment stack and operator
fusion}\label{deployment-stack-and-operator-fusion}

The raw GEMMs are at the silicon ceiling, so every remaining end-to-end
gain came from fusing the glue between GEMMs, removing the memory
round-trips a transformer block otherwise pays in eager mode.

\begin{itemize}
\tightlist
\item
  \textbf{\texttt{QuadbitLinear}} is an \texttt{nn.Linear} drop-in: a
  torch packer reproduces the kernel's exact metadata, compression, and
  scale layout (verified maxrel 0.0039), fed by a fused 128-bit NVFP4
  activation quantizer in a single CUDA pass. End-to-end it is 4.0 to
  4.2x over torch bf16 at 8192, for any token count.
\item
  \textbf{Fused SwiGLU FFN.} Quantize the activation once (shared across
  gate and up), concatenate gate and up into one GEMM, and fuse the
  epilogue (read gate and up, compute silu(gate)*up, emit the FP4-packed
  down-proj input and scales in one transposing pass). Cumulative versus
  torch bf16: unfused 2.05x, plus fused epilogue 4.45x, plus concat
  4.66x at batch 2048 (2.27x over unfused), and 1.74x to 2.95x at batch
  512. Numerically identical to the unfused path; a pure memory-traffic
  win.
\item
  \textbf{Fused RMSNorm-plus-quant} (\texttt{rmsnorm\_quant\_k}): one
  read of x replaces eager rmsnorm's read-reduce-write plus a separate
  quant pass, and it is more accurate (no bf16 round-trip before quant).
  3.7 to 4.3x over eager rmsnorm-plus-quant.
\item
  \textbf{Fused residual-add-plus-RMSNorm-plus-quant}
  (\texttt{add\_rmsnorm\_quant\_k}): the full block transition in one
  kernel. 5.3 to 5.8x over the eager sequence.
\end{itemize}

Every inter-GEMM memory round-trip in a transformer block is now a
single fused pass, at zero accuracy cost (the FP4 and prune floors are
preserved).

\textbf{On real models.} Every projection of the current frontier models
tiles onto the kernel (all hidden sizes are multiples of 256 and
head\_dim is 128): we ran the real linear shapes of Qwen3.5-397B,
GLM-5.2, MiniMax-M3, and DeepSeek-V3/R1 through the kernel, all
all-tile-ok. A full fused dense FP4 decoder block on real Qwen3-8B
weights, with no training, runs 2.16 to 2.19x over bf16 at block
reconstruction rel 0.13.

\begin{center}\rule{0.5\linewidth}{0.5pt}\end{center}

\subsection{7. Accuracy: dense FP4 is
W4A4}\label{accuracy-dense-fp4-is-w4a4}

Because the tensor core multiplies fp4 by fp4, the accuracy question is
W4A4, not W4A16. With a per-16 two-level NVFP4 recipe using amax and no
calibration data, measured through \texttt{harness/recovery\_worth.py}
on WikiText-2:

\begin{itemize}
\tightlist
\item
  Llama-3.1-8B-Instruct: 7.27 to 7.90, \textbf{+0.63 PPL}.
\item
  Meta-Llama-3-8B base: 6.20 to 6.91, \textbf{+0.71 PPL}.
\item
  Reference: vLLM native NVFP4 (modelopt-calibrated) is 7.97, +0.71 on
  the same windows.
\end{itemize}

So our own amax recipe, with no calibration, lands at or below the
calibrated reference. The earlier ``+2 PPL'' figure was a crude per-32
single-level recipe, not W4A4's real cost; the gap was block granularity
and two-level activation scaling, not calibration. The often-quoted +0.3
PPL is a weight-only W4A16 number that the hardware never runs, and we
do not headline it.

At the block-reconstruction level, two-level NVFP4 reaches rel 0.097 on
real Qwen3-8B with no training, versus 0.13 for the simpler and faster
MXFP4 (\texttt{ue8m0}) path. NVFP4 is the accuracy path; MXFP4 is the
speed path (2.15x over bf16). The single-level NVFP4 block regressed to
0.38 because the \texttt{ue4m3} scale carries a mantissa and its
per-block rounding bias accumulates across the three-matmul chain; the
two-level recipe (per-row fp32 global rescaling the per-16
\texttt{ue4m3} locals) is what fixes it, since MXFP4's power-of-two
scale has no such bias.

\begin{quote}
\textbf{In progress (reserved slot, no number claimed).} A
training-free, memory-free mixed-precision refinement that keeps the
most activation-sensitive layers at W4A16 and the rest at W4A4 is under
evaluation. Layers are ranked on decontaminated C4 and scored on
held-out WikiText-2 to avoid selection-on-test overfit, and a minimal-K
sweep deploys the smallest count of bf16-activation layers that crosses
the target, since each such layer is prefill compute paid for. No
improved PPL is reported here; the figure will land once the ablation
converges and the docs re-sync. Dense W4A4 at +0.63 PPL, zero
calibration, is the accuracy result this paper stands on.
\end{quote}

\begin{center}\rule{0.5\linewidth}{0.5pt}\end{center}

\subsection{8. Sparse recovery and the two-level deploy
fix}\label{sparse-recovery-and-the-two-level-deploy-fix}

Blackwell FP4 2:4 is pair-granular: the \texttt{mma.sp} metadata selects
at fp4-pair granularity (2 of every 4 pairs kept, not 2 of every 4
elements). This is NVIDIA's documented hardware spec, not a discovery of
ours. The consequence we do contribute as a data point: on an existing
element-2:4 checkpoint (\texttt{neuralmagic/Sparse-Llama-3.1-8B-2of4}),
pair-granular selection keeps only about 87\% of the nonzero energy, so
naive reuse gives 93.6 PPL versus 7.9 for dense FP4. Element-2:4 tooling
and pair-granular hardware are not interchangeable.

Our recovery pipeline retargets SparseGPT to pair-granular masks (keep
the two pairs of four by \texttt{w\^{}2\ /\ {[}H\^{}-1{]}\^{}2} with
Hessian error compensation), then distills from the dense teacher with
the mask frozen, then runs QAT with straight-through fake-quant of both
weights (exact kernel dequant) and activations. On TinyLlama-1.1B this
recovers to 9.60 PPL through the real sparse FP4 kernel (dense teacher
7.53), and closing the STE-versus-kernel gap (matching the QAT
activation fake-quant to the deployed quantizer bit-for-bit) was worth
about 0.43 PPL.

\textbf{The deploy gap and its fix.} Recovery trains against a two-level
fake-quant: a per-row and per-column fp32 global rescaling the per-16
\texttt{ue4m3} local scale, the same move that took dense NVFP4 from
0.38 to 0.097 block-reconstruction error. The originally deployed sparse
kernel applied only the local scale, so its outputs diverged from what
QAT trained against. We measured that deploy gap directly, running three
matmuls of one recovered Meta-Llama-3-8B checkpoint through each path:
the single-level kernel deploys at 11.89 PPL, the two-level kernel at
8.95, against a fake-quant target of 8.96. The single-level path flips
22\% of top-1 next-token predictions relative to the target; the
two-level path reproduces it (mean NLL delta 0.002, 92\% top-1
agreement). We built the missing piece, a per-row and per-column fp32
global rescale in the sparse mma epilogue ported from the working dense
two-level path. Through the two-level sparse kernel the recovered
checkpoint tracks its fake-quant within 0.02 PPL: the deploy gap is
closed. The rescale costs 2 to 10\% of throughput (worst on wide-N
shapes) and the two-level sparse kernel still beats CUTLASS 80b on every
shape (1.01 to 1.13x, all correctness-verified).

\textbf{The accuracy number, honestly deployable.} The deployed sparse
mma constrains activations to per-32 granularity on the B operand, so
the honest recipe matches the QAT activation fake-quant to per-32 rather
than the per-16 the fake-quant ceiling used. With that match,
Meta-Llama-3-8B recovers to 8.47 PPL through the two-level kernel, equal
to its fake-quant (deploy gap closed end to end). Dense W4A4
zero-training on the same model is 6.91, so deployable sparse loses to
dense by about 1.56 PPL. Extending phase-2 QAT from 2k to 5k steps did
not help (it regressed to 8.96, mild overtraining against the WikiText-2
test), so 2k is the recovery sweet spot on this corpus.

\textbf{The data lever did not pay.} To test whether the residual gap
was data-limited, we ran a full-scale diverse-corpus recovery
(decontaminated C4, phase-1 to about 196M tokens). Phase-1 flattened at
10.82 PPL on the WikiText-2 test, far above the in-distribution
WikiText-103 phase-1 (8.57), because C4 is out of distribution for that
narrow test. Diverse-corpus scale did not improve the WikiText-2 number,
so the recovery is in-distribution-bound on this metric, not simply
data-starved. This is a negative result for the diverse-data hypothesis.

The honest position: dense FP4 (+0.63, zero-training) is the accuracy
result. Sparse is a speed play (about 1.33x over our own dense) that now
deploys at its trained accuracy, since the two-level kernel closes the
deploy gap, but it still loses to dense by about 1.56 PPL on recovery
and diverse data did not close that gap on the WikiText-2 metric. Sparse
is an honest Pareto point for throughput-bound serving, not an accuracy
win.

\textbf{Can a hybrid recover most of the accuracy at some of the speed?
Training-free, no.} We asked whether selectively sparsifying only the
least-sensitive matrices (keeping the rest dense W4A4) buys a meaningful
fraction of the sparse speedup while staying close to dense accuracy. We
ranked every matrix by the PPL cost of making just it 2:4-sparse
(SparseGPT one-shot pair-granular prune, Hessian-compensated, no
training), scored on a C4 selection set disjoint from the WikiText-2
test, then sparsified least-damaging-first and traced the held-out curve
(\texttt{harness/sensitivity\_sparse.py}). The result is negative: each
matrix in isolation is nearly free (per-matrix delta-PPL ranges only
-0.13 to +0.11), but the errors compound super-linearly when stacked. On
the deployment-relevant MLP linears (dense 6.74, all-sparse one-shot
38.37), staying within +0.05 PPL allows only 3\% of MLP FLOPs to go
sparse (an estimated 1.008x), and even a +0.50 PPL budget reaches just
7\% (1.018x); including attention makes it worse (all-sparse 162.7,
attention sparsity the larger destroyer). Structurally
\texttt{down\_proj} tolerates sparsity best and \texttt{up\_proj} worst.
Because even a fully sparse model is only 1.33x over dense at the
roofline, the hybrid upside is capped there regardless; a useful hybrid
would require per-mask QAT recovery and would still be bounded by that
1.33x. The training-free free-lunch hybrid does not exist, and the
sparse Pareto point that stands (Section 5: deployed sparse beats the
best available dense FP4 on prunable weights) is not a dense-model
hybrid but a whole-matrix prunability result.

\begin{center}\rule{0.5\linewidth}{0.5pt}\end{center}

\subsection{9. End-to-end serving comparison (RTX PRO
6000)}\label{end-to-end-serving-comparison-rtx-pro-6000}

We ran real serving engines on the same card and model family, same
protocol, to place quadbit against what a practitioner would actually
deploy. All rows: Llama-3.1-8B-Instruct family, CUDA graphs on
(\texttt{enforce\_eager=False}), distinct per-request prompts (identical
prompts collapse under prefix caching and inflate prefill), prefill = B
prompts of S=2048 with 1 output token, decode = B short prompts with
GEN=128 (\texttt{ignore\_eos}), WikiText-2 PPL on the same 16x2048
windows. Memory is the loader-reported weight footprint; both engines
additionally pre-allocate a KV pool to their utilization fraction, so
total device reservation is about 86 to 89 GiB regardless of quant.

\textbf{Table A: real serving engines} (paged attention, continuous
batching, decode scheduler).

\begin{longtable}[]{@{}
  >{\raggedright\arraybackslash}p{(\linewidth - 10\tabcolsep) * \real{0.0879}}
  >{\raggedright\arraybackslash}p{(\linewidth - 10\tabcolsep) * \real{0.0769}}
  >{\raggedright\arraybackslash}p{(\linewidth - 10\tabcolsep) * \real{0.0989}}
  >{\raggedright\arraybackslash}p{(\linewidth - 10\tabcolsep) * \real{0.1099}}
  >{\raggedright\arraybackslash}p{(\linewidth - 10\tabcolsep) * \real{0.3187}}
  >{\raggedright\arraybackslash}p{(\linewidth - 10\tabcolsep) * \real{0.3077}}@{}}
\toprule\noalign{}
\begin{minipage}[b]{\linewidth}\raggedright
engine
\end{minipage} & \begin{minipage}[b]{\linewidth}\raggedright
quant
\end{minipage} & \begin{minipage}[b]{\linewidth}\raggedright
weights
\end{minipage} & \begin{minipage}[b]{\linewidth}\raggedright
WT-2 PPL
\end{minipage} & \begin{minipage}[b]{\linewidth}\raggedright
prefill tok/s (B=1/8/32/64)
\end{minipage} & \begin{minipage}[b]{\linewidth}\raggedright
decode tok/s (B=1/8/32/64)
\end{minipage} \\
\midrule\noalign{}
\endhead
\bottomrule\noalign{}
\endlastfoot
vLLM 0.21 & bf16 & 15.0 GiB & 7.267 & 10209 / 26436 / 31559 / 46880 & 88
/ 690 / 2599 / 4947 \\
vLLM 0.21 & NVFP4 (modelopt, cutlass) & 5.66 GiB & 7.974 & 13530 / 63056
/ 78028 / 116831 & 131 / 1049 / 4259 / 8465 \\
SGLang 0.5 & NVFP4 (auto → FlashInfer CUTLASS \texttt{fp4\_gemm}) &
\textasciitilde5.6 GiB & 7.97 (same ckpt) & 16829 / 59769 / 73414 /
109002 & 186 / 1491 / 5424 / 10145 \\
\end{longtable}

Native NVFP4 W4A4 is real on SM120 through both engines: vLLM binds the
\texttt{modelopt\_fp4} cutlass method (not a Marlin W4A16 fallback), and
SGLang's \texttt{auto} backend selects the FlashInfer CUTLASS
\texttt{fp4\_gemm} path (autotuned at startup, visible in its log).
Against bf16, NVFP4 is 1.7x prefill and 1.7x decode at B=64 (vLLM), for
2.6x smaller weights and +0.71 PPL. SGLang and vLLM NVFP4 trade the
lead: SGLang wins decode at every batch (10145 vs 8465 at B=64) and
low-batch prefill, vLLM wins high-batch prefill (116831 vs 109002 at
B=64).

\textbf{Table B: quadbit dense FP4 prototype} (full-forward,
PREFILL-ONLY, no decode engine).

\begin{longtable}[]{@{}
  >{\raggedright\arraybackslash}p{(\linewidth - 8\tabcolsep) * \real{0.0769}}
  >{\raggedright\arraybackslash}p{(\linewidth - 8\tabcolsep) * \real{0.0769}}
  >{\raggedright\arraybackslash}p{(\linewidth - 8\tabcolsep) * \real{0.2857}}
  >{\raggedright\arraybackslash}p{(\linewidth - 8\tabcolsep) * \real{0.2747}}
  >{\raggedright\arraybackslash}p{(\linewidth - 8\tabcolsep) * \real{0.2857}}@{}}
\toprule\noalign{}
\begin{minipage}[b]{\linewidth}\raggedright
build
\end{minipage} & \begin{minipage}[b]{\linewidth}\raggedright
quant
\end{minipage} & \begin{minipage}[b]{\linewidth}\raggedright
quantized-linear weights
\end{minipage} & \begin{minipage}[b]{\linewidth}\raggedright
through-kernel WT-2 PPL
\end{minipage} & \begin{minipage}[b]{\linewidth}\raggedright
full-model prefill tok/s
\end{minipage} \\
\midrule\noalign{}
\endhead
\bottomrule\noalign{}
\endlastfoot
quadbit dense two-level & W4A4 (per-16 ue4m3 + fp32 global) & 3.93 GiB*
& 7.899 & 7987 (B=8, S=2048) \\
\end{longtable}

*3.93 GiB counts only the transformer-block linears swapped to the FP4
kernel; embeddings, lm\_head, and attention stay bf16 and are not
counted, so this is not a full-model on-device footprint comparable to
Table A's loader figures.

Table B is deliberately not in Table A. quadbit ships a kernel plus an
\texttt{nn.Linear} drop-in, not a serving stack: no paged attention, no
continuous batching, no CUDA graphs, no decode scheduler, and attention
runs in HuggingFace eager bf16. The prefill number is
GEMM-plus-activation-quantizer bound over an eager Python forward, so it
trails the serving engines' prefill by roughly an order of magnitude and
must not be read as a serving-throughput result. What Table B does
establish is that the deployed W4A4 path is accuracy-correct end to end:
through-kernel PPL 7.90 matches vLLM's native NVFP4 7.97 to within 0.07
on the same windows, confirming quadbit's zero-calibration two-level
recipe reaches the calibrated reference's accuracy.

\textbf{Table C: quadbit inside vLLM under production CUDA graphs (one
checkpoint).} quadbit's two-level \emph{sparse} MLP is registered as a
\texttt{torch.library} custom op (\texttt{quadbit::fused\_mlp}) so
vLLM's V1 fullgraph compile and CUDA-graph capture include it (NVFP4 for
all non-MLP linears, recovered Llama-3.1-8B-Instruct; accuracy and speed
on the SAME checkpoint). The sparse path provably runs under production
graphs (\texttt{SPARSE\_CALLS=7264}, through-serving PPL
\textbf{10.2709}, not the 7.97 dense value --- the accuracy check guards
against a silent fall-back to dense NVFP4). This graph-vs-graph
comparison is the production-representative serving result (util 0.8,
WT-2 15x2048); see \texttt{docs/graph\_serving\_result.md}.

\begin{longtable}[]{@{}
  >{\raggedright\arraybackslash}p{(\linewidth - 6\tabcolsep) * \real{0.0941}}
  >{\raggedright\arraybackslash}p{(\linewidth - 6\tabcolsep) * \real{0.3647}}
  >{\raggedright\arraybackslash}p{(\linewidth - 6\tabcolsep) * \real{0.5059}}
  >{\raggedright\arraybackslash}p{(\linewidth - 6\tabcolsep) * \real{0.0353}}@{}}
\toprule\noalign{}
\begin{minipage}[b]{\linewidth}\raggedright
metric
\end{minipage} & \begin{minipage}[b]{\linewidth}\raggedright
vLLM NVFP4 (graph, production)
\end{minipage} & \begin{minipage}[b]{\linewidth}\raggedright
quadbit sparse MLP + split-K down (graph)
\end{minipage} & \begin{minipage}[b]{\linewidth}\raggedright
Δ
\end{minipage} \\
\midrule\noalign{}
\endhead
\bottomrule\noalign{}
\endlastfoot
WT-2 PPL & 7.97 & 10.2709 & +2.30 \\
prefill B=8/32/64 (tok/s) & 66469/80825/119083 & 62914/77605/115069 &
\textbf{−5.3\% / −4.0\% / −3.4\%} \\
decode B=8/32/64 (tok/s) & 1046/4237/8384 & \textbf{1147/4543/8567} &
\textbf{+9.7\% / +7.2\% / +2.2\%} \\
\end{longtable}

\textbf{quadbit beats production dense NVFP4 on decode} --- the
latency-critical, memory-bound serving regime --- at every batch, with
correct sparse accuracy. Prefill trails \textasciitilde3-5\% (it uses
the plain down, which was never underfilled). The decode win comes from
a split-K down projection; the earlier −6 to −12\% decode loss was a
diagnosed underfill, now fixed.

\emph{Decode diagnosis and fix (\texttt{-\/-mode\ profile\_decode}, CUDA
events, µs/layer; kernel time is graph-invariant):} activation-quant 15,
gate\_up sparse GEMM 52 (\textbf{112 CTAs}), fused SwiGLU 30,
\textbf{plain down sparse GEMM 109 (16 CTAs)}. \texttt{matmul\_sp}
launches \texttt{grid=(N/BN,\ M/2BM)}, BM=BN=128; at decode (tp=128)
occupancy is set by the output dim, so gate\_up (28672) fills 112 CTAs
but \textbf{down (4096) underfilled at 16 CTAs on \textasciitilde188 SMs
(\textasciitilde8\% of the GPU)}. The fix is a \textbf{split-K down
kernel} (\texttt{matmul\_sp\_sk}): a \texttt{gridDim.z} K-split raises
the down to 16 × 8 = \textbf{128 CTAs}, an f32 workspace accumulates the
partial sums, and a \texttt{cvt\_sp\_2lvl\_t} pass applies the two-level
global scales post-reduction and transposes to
\texttt{{[}tok,\ out\_f{]}}. Down drops \textbf{109 → 56.5 µs (1.94×)}
at split=8 (cos 1.0000, relL2 ≈ 1e-5 vs the plain two-level down), and a
serving sweep confirms split=8 beats 4 and 16 end-to-end. The op
dispatches to \texttt{fused\_mlp\_2lvl\_skdown} at decode (tp ≤ 128) and
the plain \texttt{fused\_mlp\_2lvl} at prefill.

\emph{Eager-vs-eager (diagnostic ablation, not the production serving
headline).} The eager path's optimization story explains the kernel
work: (1) a zero-copy transposed mma epilogue
(\texttt{sparse\_fp4\_mm\_2lvl\_t}), bitwise- equal to \texttt{.t()};
(2) a two-level fused SwiGLU (11.13 single-level → 8.95 on base); (3) a
single no-sync \texttt{fused\_mlp\_2lvl} removing \textasciitilde64
\texttt{cudaDeviceSynchronize}/forward. These took the \emph{eager} path
to +3.7/+5.5/+5.6\% prefill and +23\% decode vs \emph{eager} NVFP4
(Table D-eager / \texttt{docs/frozen\_serving\_result.md}). That win is
\textbf{launch-overhead only} and does not survive once both paths are
CUDA-graphed --- hence Table C above, not the eager numbers, is the
serving result.

\textbf{Table C decode/prefill split hides the request-level answer, so
we measured the crossover directly.} Table C reports prefill and decode
as separate throughputs; a real request pays both, so which path wins
\emph{end-to-end} depends on the decode fraction. We swept a batch x
prompt-length x generation-length request matrix, both paths CUDA-graph
captured, scoring total request latency per cell (TTFT from
\texttt{generate(max\_tokens=1)}, total from
\texttt{generate(max\_tokens=G,\ ignore\_eos=True)}). Every
decode-throughput number in this paper is \textbf{decode-only}, across
the Llama crossover here, the DeepSeek serving sweep (Section 10), and
the GLM serving rows (Section 10.1, \texttt{docs/glm\_results.md}). The
\texttt{generate(max\_tokens=1)} (prefill/TTFT) wall is subtracted from
the \texttt{max\_tokens=G} wall so the reported rate reflects the decode
steps alone, never prompt+generation throughput. \textbf{Prefix caching
must be OFF}: with it on, vLLM V1 reuses the TTFT call's prompt KV in
the total call, skips the real prefill, and hides sparse's prefill
deficit --- a first pass with caching on spuriously showed sparse
winning all 112 cells. With it off (each request pays a real prefill),
quadbit's sparse split-K FP4 MLP \textbf{wins end-to-end total request
latency outright in 81 of 112 regimes and ties 2 more (83/112 where it
is at least as fast)} vs production dense NVFP4. B=1 single-stream wins
\emph{every} regime (+3.5\% to +11.6\%); sparse wins at any batch once
generation clears a batch/prompt-dependent boundary; NVFP4 keeps only
the prefill-bound corner (high batch x long prompt x short generation,
where it wins by \textless=3\%).

\emph{Crossover boundary --- min generation length for sparse to win
total latency:}

\begin{longtable}[]{@{}lllll@{}}
\toprule\noalign{}
B & prompt=128 & 512 & 2048 & 8192 \\
\midrule\noalign{}
\endhead
\bottomrule\noalign{}
\endlastfoot
1 & 16 & 16 & 16 & 16 \\
8 & 16 & 16 & 32 & 128 \\
32 & 16 & 32 & 128 & 128 \\
64 & 16 & never (tie at 256) & 1024 & never (\textless=1024) \\
\end{longtable}

The boundary rises with batch and prompt length: as the workload becomes
more prefill-bound, sparse needs a longer generation to amortize its
prefill deficit. Accuracy is a constant +2.3 PPL (10.27 vs 7.97) across
the whole map, so the crossover is purely a speed map. \textbf{The
serving claim: for interactive / low-batch and long-generation regimes,
sparse FP4 wins end-to-end}; the batch-prefill corner stays NVFP4's.
Full matrix in \texttt{docs/crossover\_result.md}.

\textbf{Verification / multi-token decode does \emph{not} favor sparse
(refuted).} Speculative/verification decoding processes k candidate
tokens per sequence, so the MLP sees effective M = B\emph{k rows per
decode step; the natural hypothesis is that larger M favors sparse
tensor-core work. We measured it and it is false: the sparse/NVFP4
decode-throughput margin }shrinks* with M (M=1: +13\%, M=64: +2\%, noisy
and NVFP4-favorable beyond) and never expands. The split-K decode win is
a \textbf{small-M GPU-underfill fix} that fades as M grows and NVFP4's
own dense GEMM fills the machine. So sparse FP4 is best for
\textbf{low-M latency-sensitive decode}, not throughput-oriented
multi-token verification --- consistent with the crossover map (sparse
dominates B=1-8; the batch-heavy corner is NVFP4's).

\textbf{Attacking the +2.3 PPL tax by reverse densification: no free
Pareto point (Section 8 confirmed at serving scale).} To close the
constant tax we reverted selected MLP projections from recovered-sparse
back to stock dense NVFP4, measured through the serving path. All-sparse
10.256 PPL, all-dense 7.974. Densifying \texttt{down\_proj} recovers
\textasciitilde0 PPL (down-sparsity is accuracy-free \emph{and} is
exactly where the split-K decode win lives), while \texttt{gate\_up}
carries the recoverable tax (all \texttt{gate\_up} dense → 9.750, -0.51,
late layers cost most). But densifying \texttt{gate\_up} \emph{hurts}
decode by 7 to 9\%, because sparse \texttt{gate\_up} was already the
fast component. Reverse densification therefore trades speed for
accuracy roughly 1:1 with no free knee; closing the tax needs QAT repair
of the \texttt{gate\_up}-dense/\texttt{down}-sparse hybrid (9.750), not
placement alone. This is consistent with the training-free
hybrid-placement negative result (Section 8). Full Pareto in
\texttt{docs/accuracy\_pareto.md}.

\textbf{Phase-adaptive same-weight execution (Track 4C): built and
refuted.} A tempting way to erase the prefill-bound corner NVFP4 keeps
is to run the \emph{same} recovered pruned weights in two layouts by
effective token count: prefill (large M) through a production dense
NVFP4 GEMM (FlashInfer \texttt{cutlass}, weights materialized dense with
zeros in the pruned slots), decode (small M) through the sparse split-K
path. We built it. The semantics hold: dense NVFP4 over the recovered
weights scores 10.30 PPL through serving, equal to all-sparse 10.27
within numerical drift, so the two layouts are interchangeable and the
phase boundary is seamless. But the row loses: 39 wins and 66 losses of
105 crossover cells versus NVFP4, against all-sparse's 81 wins and 29
losses, and it flips none of the cells all-sparse already lost. The root
cause is measured (\texttt{-\/-mode\ phase\_bench}, microseconds per MLP
layer at prefill): the hand-rolled dense path (\texttt{nvfp4\_quantize}
plus \texttt{mm\_fp4} plus SwiGLU plus \texttt{nvfp4\_quantize} plus
\texttt{mm\_fp4}) runs about 2x native NVFP4 because its activation
quant is unfused (the FlashInfer \texttt{nvfp4\_quantize} of the down
input alone is 517 microseconds at M=2048, more than both GEMMs
combined), whereas vLLM fuses that quant into the norm and the SwiGLU
through compiled passes an opaque custom op cannot reach. Decisive: the
sparse fused MLP is already about 7 to 10 percent faster per layer than
native NVFP4 (618 versus 661 microseconds at M=2048), so there is no
faster dense MLP to swap in, and the corner all-sparse loses is
attention and Amdahl bound, not MLP bound. Track 4C is a documented dead
end, not a win. Full analysis in \texttt{docs/crossover\_result.md}
Section 4C.

\textbf{The open axis: repairing the +2.3 PPL tax.} Every serving result
above carries the same constant accuracy tax, +2.3 PPL (10.27 sparse
versus 7.97 dense NVFP4), and the crossover is purely a speed map
because that tax does not move across the request matrix. Two
training-free levers to close it are now measured and negative: reverse
densification (this section) trades speed for accuracy about 1:1, and
training-free hybrid placement (Section 8) buys almost no sparse FLOPs
before errors compound. The remaining paper-upgrade axis is therefore
accuracy repair that spends training. A tournament of four approaches
ran on the recovered-Instruct all-sparse checkpoint: (1) zero-runtime
calibration of the deployed sparse recipe, (2) low-rank residual
adapters over the pruned weights, (3) activation-aware (Wanda-pair) mask
repair, and (4) knowledge distillation from the dense teacher that
retrains the sparse weights and the per-output down scale.

\textbf{Result: distillation reduces the perplexity tax but does not
recover downstream capability.} Only the distillation family moved the
metric. The best variant (KL-light/CE-heavy) reaches through-kernel PPL
\textbf{8.86} (serving PPL \textbf{9.10}, from the original
\textbf{10.27}), clearing the target and keeping the entire split-K
decode win intact (decode +9.7/+7.2/+2.2\% at B=8/32/64, the banked
figure; serving speed is weight-value independent, so the 81/112
crossover carries over unchanged). But the downstream-task check tells
the real story: on ARC-Challenge, HellaSwag, PIQA, and Winogrande the
repaired checkpoint is essentially unchanged from the un-repaired
all-sparse model (about +0.005 accuracy on average), and the lowest-PPL
variant even regressed on ARC-Challenge (0.356 to 0.348). The 2:4
sparsity removes roughly 20 points of ARC-C/HellaSwag accuracy versus
dense, and WikiText knowledge distillation recovers almost none of it.
The CE-heavy WikiText PPL win is largely domain overfitting. The other
three families were negative: zero-runtime output calibration is neutral
to harmful (per-channel affine 12.97, because an output rescale cannot
fix a representational loss), low-rank adapters on frozen sparse weights
stayed flat at about 10.0, and Wanda-pair masks with truncated QAT were
worse than the SparseGPT baseline (13.06).

\begin{quote}
\textbf{Accuracy-repair result: PPL repaired, capability not.}
Distillation reduces the sparse serving tax from 10.27 to 9.10 PPL and
preserves every serving win, but it does not close the
downstream-capability gap that 2:4 sparsity opens (ARC-C/HellaSwag stay
\textasciitilde20 points below dense). Distillation on a narrow text
corpus buys perplexity, not capability. The honest remaining frontier is
sparse capability recovery (broad/larger-scale distillation data, or
rethinking the prune target), not serving plumbing. See
\texttt{docs/crossover\_result.md} and \texttt{harness/repair.py}.
\end{quote}

\begin{center}\rule{0.5\linewidth}{0.5pt}\end{center}

\subsection{10. Cross-architecture sparse-policy transfer
(DeepSeek-V4-Flash and
GLM-5.2)}\label{cross-architecture-sparse-policy-transfer-deepseek-v4-flash-and-glm-5.2}

Sections 4 to 9 establish sparse FP4 on a single dense model
(Llama-3-8B) on one GPU. The obvious reviewer questions are whether the
approach (i) transfers to a large model, (ii) transfers to a
Mixture-of-Experts architecture, and (iii) works across multiple GPUs.
We answer all three on \textbf{DeepSeek-V4-Flash} (284B total / 13B
active, 256 routed + 1 shared expert per MoE layer, top-6,
\texttt{moe\_intermediate} 2048, hidden 4096, 43 layers, MLA +
sparse-index attention). This is a strictly harder setting than
Llama-3-8B: the MLP FLOPs live in hundreds of routed experts behind a
router, not a single \texttt{LlamaMLP}, so the dense-MLP operator of
Sections 6/9 does not attach.

\textbf{The experts are MXFP4, not the config's FP8.} The checkpoint's
\texttt{quantization\_config} advertises FP8 block-128, but that
describes the dense/attention linears; the experts
(\texttt{expert\_dtype:\ fp4}) are \textbf{MXFP4} -- E2M1 codes packed
two per int8 byte with an e8m0 (power-of-two) scale per 32-element
block. We decode to bf16 (low nibble -\textgreater{} even element,
matching the OCP/transformers convention) and verify the decode is
\textbf{value-exact (100\% round-trip)} on real weights, then prune 2:4
by magnitude and re-quantize into the two-level NVFP4 layout of Section
3. The re-quantized experts match the through-kernel path at cos 0.9999.

\textbf{A segmented routed-row sparse GEMM makes MoE graph-capturable.}
A per-expert Python loop is correct but not CUDA-graph-capturable
(data-dependent launch count) -- fatal for the graph-mode serving of
Section 9. We instead stack the local experts' packed weights into one
tensor \texttt{{[}E*Mpe,\ ...{]}}, sort the routed rows by expert, pad
each segment to the 128-row column tile, and pass a per-column-block
expert id. A single kernel (\texttt{matmul\_sp\_moe}, one launch, fixed
grid) reads its expert id and offsets only the weight-side indices
(TMA-A coordinate, local scales, 2:4 metadata, per-row global scale) by
\texttt{expert*Mpe}, keeping output-side indices local. The scheduling
unit is a routed row, not an expert. Validated bit-exact against the
per-expert kernel (\textbf{cos 1.000000}) across uniform, all-to-one,
one-per-expert, and imbalanced routing at both tiny and real DeepSeek
expert shapes; \textbf{CUDA-graph capture and replay reproduce it at cos
1.000000}. On one real layer (256 experts, real gate routing) the
segmented operator matches the per-expert kernel at cos 1.000000 with
zero non-finites.

\textbf{Coverage.} quadbit sparsifies every routed and shared expert
MLP; the router, MLA attention, embeddings and lm\_head stay dense
NVFP4. That is \textbf{\textasciitilde91\% of model parameters} and
\textbf{\textasciitilde80\% of active linear FLOPs per token} (top-6/256
experts + shared). In dense Llama-3-8B the MLP is \textasciitilde66\% of
parameters, so the sparse coverage is \emph{higher} for the large MoE,
not lower -- the result transfers, and transfers to more of the model.

\textbf{Distributed (expert-parallel) scaling.} Sharding the 256 experts
across ranks (expert parallelism, the natural choice on RTX PRO 6000,
which has no NVLink) and combining per-rank contributions with an
all-reduce, the segmented kernel scales near-linearly: \textbf{2.17x on
2 GPUs and 4.21x on 4 GPUs} of expert-kernel time, routing imbalance
1.04, with identical output checksum at every world size (correctness
preserved). On PCIe the all-reduce communication is 0.32 to 0.45 ms
against 1.3 to 2.5 ms of expert compute -- non-trivial but not dominant
at this scale; on a no-NVLink fabric this communication share is the
honest external-validity caveat, and it grows with world size.

\textbf{Accuracy.} The per-expert-output 2:4-FP4 error on real weights
(random activations, a worst case) is cos \textasciitilde0.70 --
consistent with the single-model finding that 2:4 sparsity, not FP4, is
the accuracy cost. Naively sparsifying every expert compounds this over
43 layers to incoherence. Section 10.1 shows this is \textbf{recoverable
training-free} by placing the sparsity structurally (which projections,
which layers, which routed slots), not by repairing weights -- the
per-expert repair we tried fails, but the capability-preserving
operating points do not need it.

\textbf{Serving integration status (an honest ecosystem boundary).} In
vLLM 0.24 the DeepSeek-V4-Flash NVFP4 checkpoint downloads, and its
weights load across 2x RTX PRO 6000; vLLM selects the FlashInfer CUTLASS
NVFP4 MoE backend and the FP8 Lightning-Indexer attention. But the
model's \textbf{FP8 (ue8m0 W8A8) attention/dense GEMMs have no working
kernel on SM120}: with DeepGEMM enabled the ue8m0 scale-factor transform
asserts (\texttt{Unknown\ SF\ transformation}), and with DeepGEMM
disabled the fallback CUTLASS c3x \texttt{scaled\_mm} has no SM120
dispatch (\texttt{dispatch\_scaled\_mm}). Both paths are
Hopper-targeted, so the model cannot complete even vLLM's init profiling
forward on consumer Blackwell -- a limitation of the FP8 attention path
in today's serving stack, \textbf{independent of quadbit} (the MoE path
that quadbit replaces is the one backend that did initialize). We
therefore report the quadbit sparse MoE result where it is measured
cleanly: the operator is validated standalone (segmented kernel
bit-exact vs per-expert; CUDA-graph capturable; correct on a real
256-expert layer) and \textbf{distributed} (expert-parallel, 4.21x
kernel scaling on 4 GPUs, correctness preserved), with the coverage and
accuracy accounting above. The staged-\texttt{.so} injection path
(compile the sparse mma under CUDA \textless= 12.8, ctypes-load under
the CUDA-13 serve image) and the FusedMoE hook are implemented
(\texttt{harness/serve\_dsv4.py}). \textbf{Section 10.1 removes this
gate:} we supply the missing SM120 paths ourselves through a vLLM plugin
and serve the model end-to-end.

\subsubsection{10.1 Overturning the ecosystem boundary, and
training-free capability
preservation}\label{overturning-the-ecosystem-boundary-and-training-free-capability-preservation}

The ``future work gated on ecosystem FP8 SM120 support'' caveat above is
resolved by building the missing kernels rather than waiting for them. A
vLLM \texttt{general\_plugins} plugin (installed in every spawned
worker) supplies each absent SM120 path: block-FP8 dense/attention
linears dequantised to bf16 at load; a hand-reimplemented MLA
\texttt{o\_proj} (inverse-interleaved RoPE + block-dequant einsum)
replacing the Hopper-only \texttt{deep\_gemm\_fp8\_o\_proj}; the
DeepSeek sparse-attention Lightning-Indexer logits reimplemented in bf16
(the FlashInfer sparse-MLA core \emph{is} SM120-supported, only the
indexer needed owning); and a pure-torch override of the
cooperative-cluster top-k kernels that fail to launch on SM120. With
these, \textbf{DeepSeek-V4-Flash-NVFP4 generates coherent text
end-to-end in vLLM on 2x RTX PRO 6000}, and the quadbit 2:4-sparse-FP4
experts run in the live serving loop (load drops from \textasciitilde84
to \textasciitilde56 GiB/GPU as raw NVFP4 is freed after packing --
proof the sparse path executed, not a silent fallback).

\textbf{Training-free capability preservation.} Downstream evaluation
(400 items/task over ARC-C, HellaSwag, Winogrande, MMLU-5; dense AVG
.7383) shows the accuracy cost is governed by \emph{where} the sparsity
is placed, and recovers with no weight training:

\begin{itemize}
\tightlist
\item
  \textbf{Projection anchoring.} The downstream tax lives in the gate/up
  projection; the down projection is nearly free. Sparsifying only down
  projections in the later 49\% of MoE layers (\textbf{c\_down49}) holds
  \textbf{.7354, -0.29pt} from dense on 2 GPUs, serving path and memory
  unchanged; the gate/up-only control at the same coverage falls to
  .7056 (-3.27pt). Down-only clears the -2pt bar up to 60\% of layers;
  the cliff is at 65\%.
\item
  \textbf{Route-slot.} Keeping the top-2 highest-weight routed slots per
  token dense and 2:4-sparsifying the low-weight tail (\textbf{D2})
  reaches \textbf{\textasciitilde33\% active sparse expert-FLOP at
  -0.79pt} (.7304), double c\_down49's FLOP share -- the dominant
  experts carry capability, the tail is nearly free. This needs dense
  and sparse weights co-resident for the same experts (dual residency),
  so it runs at 4 GPUs.
\item
  A per-expert layerwise repair (local KD on each expert's surviving 2:4
  weights) \textbf{fails} (-7pt): it trades the checkpoint's global
  weight consistency for local optima that compound. Structural
  placement succeeds where weight repair does not. Magnitude/Wanda-alone
  and all-expert sparsity also fail (all \textless= -4pt). See
  \texttt{docs/deepseek\_final\_table.md} and Figure
  \texttt{fig\_ds\_pareto} / \texttt{fig\_ds\_designspace}.
\end{itemize}

\textbf{Transfer to GLM-5.2.} To test that these are model-general
structural rules and not DeepSeek idiosyncrasies, we repeat them on
\textbf{GLM-5.2-NVFP4} (\texttt{glm\_moe\_dsa}: 78 layers / 75 MoE,
256+1 experts, top-8, Deep Sparse Attention + MLA, 432.9 GiB). GLM loads
and generates coherently on \textbf{8x RTX PRO 6000} (expert-parallel);
its DSA runs natively on SM120 (vLLM selects
\texttt{FLASHINFER\_MLA\_SPARSE\_SM120}), no fallback. Anchoring MoE
layers 0-37 dense and sparsifying 38-74 (49.3\%), held-out PPL (dense
3.171) moves by \textbf{+0.209 for down-only vs +0.432 for gate/up} --
the same ``down safe, gate/up expensive'' mechanism at roughly half the
tax -- while route-slot D2 (top-2 dense, tail 2:4) gives the
\textbf{best quality and the highest sparse-FLOP together (+0.065 PPL,
\textasciitilde37\% FLOP)}, mirroring DeepSeek's D2. Figure
\texttt{fig\_glm\_transfer} places the two models side by side.
\textbf{Route-slot D2 is the headline GLM policy}, and to check that its
small PPL cost is not hiding a downstream collapse we run the
tokenizer-agnostic MC harness on GLM (the same
ARC-C/HellaSwag/Winogrande/MMLU-5 suite as DeepSeek, PIQA excluded
because \texttt{ybisk/piqa} is not loadable on the serve image, so both
are 4-task; \texttt{limit=200}):

\begin{longtable}[]{@{}lllllll@{}}
\toprule\noalign{}
GLM policy & ARC-C & HellaSwag & Winogrande & MMLU-5 & AVG & PPL \\
\midrule\noalign{}
\endhead
\bottomrule\noalign{}
\endlastfoot
dense (ref) & .655 & .780 & .750 & .856 & \textbf{.7603} & 3.171 \\
\textbf{route-slot D2} & .650 & .780 & .725 & .848 & \textbf{.7508} &
3.216 \\
\end{longtable}

D2 holds \textbf{within 0.95 pt AVG of dense with no task collapsing}
(HellaSwag exactly flat; ARC-C and MMLU-5 within the n=200 /
n=50-per-subject band; Winogrande's -2.5 pt is the largest single move
and near its per-200 noise), so the small PPL gap does not mask a
downstream regression. One honest limit remains on quality: this GLM
downstream evidence is a small 4-task smoke suite on the D2 policy only
(full benchmarks, and downstream numbers for the down-only/gate-up rows,
are still unmeasured on GLM, so the most complete capability accounting
of record remains DeepSeek's).

\textbf{Graph-enabled and dense-anchor delegation (P4 + C1).} The
PPL/downstream rows above are the deployed quality reference; the
serving path itself is now graph-enabled and no longer dense-loop-bound.
P4 replaced the plugin's expert-parallel host-sync
(\texttt{torch.unique(...).tolist()}, illegal under CUDA-graph stream
capture) with a fixed-capacity device-routing path, so the deployed
sparse policies CUDA-graph capture on SM120. C1 then removed the
remaining dense-anchor decode bottleneck: the anchored/grouped
projection previously ran a dequant-to-bf16 loop over all local experts,
and we replaced it with FlashInfer's native grouped NVFP4 GEMM
(\texttt{group\_gemm\_nvfp4\_nt\_groupwise}, opt-in
\texttt{QB\_DENSE\_BACKEND=native\_nvfp4}), with no custom dense
grouped-GEMM. In standalone A/B at DeepSeek shapes the native path
matches the dequant loop (cos 0.991 vs bf16, no non-finite),
graph-captures with bit-identical replay, and is \textasciitilde18-25x
faster. In serving, the native-delegate captured paths decode faster
than eager at matched quality:

\begin{longtable}[]{@{}
  >{\raggedright\arraybackslash}p{(\linewidth - 8\tabcolsep) * \real{0.1667}}
  >{\raggedleft\arraybackslash}p{(\linewidth - 8\tabcolsep) * \real{0.2222}}
  >{\raggedleft\arraybackslash}p{(\linewidth - 8\tabcolsep) * \real{0.2222}}
  >{\raggedleft\arraybackslash}p{(\linewidth - 8\tabcolsep) * \real{0.2222}}
  >{\raggedright\arraybackslash}p{(\linewidth - 8\tabcolsep) * \real{0.1667}}@{}}
\toprule\noalign{}
\begin{minipage}[b]{\linewidth}\raggedright
config
\end{minipage} & \begin{minipage}[b]{\linewidth}\raggedleft
GPUs
\end{minipage} & \begin{minipage}[b]{\linewidth}\raggedleft
PPL
\end{minipage} & \begin{minipage}[b]{\linewidth}\raggedleft
decode tok/s
\end{minipage} & \begin{minipage}[b]{\linewidth}\raggedright
graph
\end{minipage} \\
\midrule\noalign{}
\endhead
\bottomrule\noalign{}
\endlastfoot
DeepSeek D2 dequant, captured (base) & 4 & 3.9746 & 0.514 & FULL \\
DeepSeek D2 native, eager & 4 & 4.0483 & 1.637 & none \\
\textbf{DeepSeek D2 native, captured} & 4 & \textbf{4.0112} &
\textbf{5.820} & FULL \\
\textbf{GLM route-slot D2 native, captured} & 8 & \textbf{4.0705} &
\textbf{5.296} & PIECEWISE 3/3 + FULL 2/2 \\
\end{longtable}

Native-captured DeepSeek-D2 is \textbf{11.3x the dequant-captured
baseline} (same harness) and \textbf{1.44x the frozen-eager 4.04 tok/s}
at matched PPL; the speedup decomposes as native backend 3.2x (the
grouped GEMM is far cheaper than the dequant loop) times capture 3.6x
(launch-overhead removal, which only pays once each step is cheap). GLM
route-slot D2 native-captured decodes at \textbf{5.296 tok/s = 2.5x the
eager reference 2.10}, with DSA
\texttt{sparse\_mla\_sm120\_decode\_dsv3\_2} native and pool 1.21
GiB/GPU. These are same-model/same-policy speed/quality/memory Pareto
results (native delegate vs our own dequant-loop and eager paths), not a
production-wide decode-speed claim over other serving stacks. Two honest
limits remain on the serving path: the native delegate depends on
FlashInfer availability and its swizzled NVFP4 scale layout, and the GLM
graph rows are validated on a short held-out passage (the dense-baseline
3.171 uses a different 114-token policy-sweep passage, so only
within-protocol capture-neutrality comparisons are valid). Full result:
\texttt{docs/c1/verdict.md}; logs \texttt{docs/audit/logs/c1\_*.log}.

\textbf{C2 SOTA board: the dense NVFP4 fused MoE is the SM120 MoE decode
SOTA, not the sparse policy.} To place the C1 result honestly, we ran a
direct same-harness board (same mito80 passage, decode-only formula,
captured graph mode, memory accounting) against the strongest
dense/NVFP4 baseline that runs on SM120: vLLM's native
FlashInfer-CUTLASS fused NVFP4 MoE (\texttt{QB\_MOE=off}, with the
plugin's attention/DSA unblock so the model loads at all; vanilla vLLM
cannot init this model on SM120, and SGLang has no SM120 path for it).

\begin{longtable}[]{@{}
  >{\raggedright\arraybackslash}p{(\linewidth - 6\tabcolsep) * \real{0.2143}}
  >{\raggedleft\arraybackslash}p{(\linewidth - 6\tabcolsep) * \real{0.2857}}
  >{\raggedleft\arraybackslash}p{(\linewidth - 6\tabcolsep) * \real{0.2857}}
  >{\raggedright\arraybackslash}p{(\linewidth - 6\tabcolsep) * \real{0.2143}}@{}}
\toprule\noalign{}
\begin{minipage}[b]{\linewidth}\raggedright
model (captured)
\end{minipage} & \begin{minipage}[b]{\linewidth}\raggedleft
dense NVFP4 fused
\end{minipage} & \begin{minipage}[b]{\linewidth}\raggedleft
quadbit D2 native
\end{minipage} & \begin{minipage}[b]{\linewidth}\raggedright
dense advantage
\end{minipage} \\
\midrule\noalign{}
\endhead
\bottomrule\noalign{}
\endlastfoot
DeepSeek, 4 GPU & 48.248 tok/s, 40.83 GiB/GPU & 5.972 tok/s, 51.7
GiB/GPU & 8.1x decode, less memory \\
GLM, 8 GPU & 33.810 tok/s, 54.62 GiB/GPU & 5.367 tok/s, 68.98 GiB/GPU &
6.3x decode, less memory \\
\end{longtable}

The dense fused MoE wins decode (6-8x) and weight memory (the sparse
policy's dual residency costs +26-27\% weight memory and, on GLM, 62\%
less KV). C1's 11.3x was measured against the crippled dequant loop
(0.514 tok/s), not this dense path. So the sparse route-slot policy is
\textbf{not} a decode-speed or decode-memory SOTA; its contribution is
the only deployed 2:4-sparse FP4 MoE, training-free
capability-preserving structural sparsity that transfers across
architectures and graph-captures with native DSA (downstream within
\textasciitilde1 pt), and the prefill/large-M kernel Pareto of Section 5
(a kernel microbench, not re-measured in serving here). The remaining
MoE decode gap is a kernel problem, a fused sparse grouped decode GEMM
at tiny M, identified but not built. Board:
\href{c2/verdict.md}{docs/c2/verdict.md}; logs
\texttt{docs/audit/logs/c2\_*.log}.

\subsubsection{10.2 Layerwise QAT does not clear the gate/up
floor}\label{layerwise-qat-does-not-clear-the-gateup-floor}

The gate/up tax of Section 10.1 raises the obvious question: can
quantization-aware training recover what training-free placement avoids?
We ran the whole tractable family of layerwise-repair objectives to the
ground on a widened 8-task downstream battery (ARC-C, ARC-E, HellaSwag,
PIQA, OBQA, BoolQ, Winogrande, MMLU-5; \texttt{limit=400}; dense NVFP4
AVG \textbf{.7548}), and every one is a negative:

\begin{longtable}[]{@{}
  >{\raggedright\arraybackslash}p{(\linewidth - 4\tabcolsep) * \real{0.3333}}
  >{\raggedright\arraybackslash}p{(\linewidth - 4\tabcolsep) * \real{0.3333}}
  >{\raggedright\arraybackslash}p{(\linewidth - 4\tabcolsep) * \real{0.3333}}@{}}
\toprule\noalign{}
\begin{minipage}[b]{\linewidth}\raggedright
gate/up recovery objective (DeepSeek-V4-Flash-NVFP4)
\end{minipage} & \begin{minipage}[b]{\linewidth}\raggedright
AVG-8
\end{minipage} & \begin{minipage}[b]{\linewidth}\raggedright
vs one-shot
\end{minipage} \\
\midrule\noalign{}
\endhead
\bottomrule\noalign{}
\endlastfoot
one-shot sparse gate/up (no training) & .7363 & --- \\
per-expert STE-QAT (each expert fit to its own dense output through the
kernel) & .7332 & −0.31pt \\
\textbf{global router-weighted MoE QAT} (top-k combine fit to the dense
routed aggregate) & \textbf{.7259} & \textbf{−1.04pt} \\
\end{longtable}

The global objective is the one the per-expert fit discards routing
weights for: it fits the router-weighted top-k combine to the dense
routed aggregate the next layer actually sees, trained by Gauss-Seidel
coordinate descent (one expert resident at a time; a joint 256-expert
backward OOMs, and a Jacobi all-vs-round-start update diverges). It
produced the \textbf{tightest reconstruction of any variant}
(routed-aggregate relative error driven monotonically to 0.004-0.017,
roughly 4-10x tighter than per-expert's 0.02-0.07) yet the
\textbf{worst} downstream of the three, moving accuracy the wrong way.
The same negative holds on a single dense model: full-stack
through-kernel QAT on Llama-3.1-8B (the Gap-C harness) lands
\textbf{0.3967 \textless{} the one-shot 0.4333}, below the very baseline
it was meant to beat.

The lesson is consistent across all four: \textbf{local reconstruction
error does not predict downstream capability, and driving it tighter
makes the mismatch worse} -- matching a calibration-trajectory target
overfits layer-local statistics that do not transfer to held-out
multiple-choice accuracy. No objective in the layerwise-repair family
(own-output per-expert, routed-aggregate global, dense- or
sparse-trajectory, single dense model) clears the 2:4 gate/up floor; the
floor is a real capability property, not a reconstruction-objective
artifact. You avoid the tax by anchoring gate/up dense (down-only or
route-slot, Section 10.1), you do not repair it. True end-to-end QAT
(attention and MLP jointly trainable \emph{through} the sparse kernel)
is \textbf{not} refuted -- it stays infra-blocked (no differentiable
DeepSeek/GLM load exists; the serving path is inference-only) -- but
this closes every layerwise objective the inference-only path can
express. Verdicts:
\texttt{docs/qat/\{design,gateup\_moe\_verdict,global\_moe\_verdict\}.md}.

\subsubsection{10.3 Transfer to gpt-oss-20b: full end-to-end serving and
honest
throughput}\label{transfer-to-gpt-oss-20b-full-end-to-end-serving-and-honest-throughput}

A third architecture, and the one we serve most completely.
\textbf{gpt-oss-20b} (24 layers, hidden 2880, 32 experts, top-4, expert
intermediate 2880) quantizes \textbf{only the expert MLPs to MXFP4} --
self-attention, router, embeddings, and lm\_head stay dense -- which is
\emph{exactly} quadbit's sparsify-experts /
keep-attention-router-embed-dense policy, arrived at independently by
the checkpoint. Its fused 3D expert tensors store gate/up interleaved
and use a clamped SwiGLU (alpha 1.702, limit 7.0). We decode the MXFP4
experts \textbf{value-exact (100\% round-trip)}, de-interleave gate/up,
zero-pad 2880 -\textgreater{} 3072x2944 for kernel tile alignment
(\textbf{padding-exact}: padded-real experts reach cos 0.987-0.990 vs
fake-quant, \emph{beating} an aligned-no-pad control at 0.978, so the
pad adds no error), and pack into the two-level NVFP4 layout of Section
3. All offline gates pass on real weights (full-expert forward cos 0.99
vs fake-quant; the actual segmented serving kernel with top-4 routing
cos 0.972 vs dense).

\textbf{Served end-to-end in vLLM 0.24}, by patching the native
\texttt{GptOssMxfp4MoEMethod}: gpt-oss-20b generates coherent text with
the 2:4-sparse experts live, raw NVFP4 freed after packing (proof the
sparse path executes, not a silent fallback). Unlike the
ecosystem-blocked DeepSeek path (Section 10), gpt-oss serves cleanly on
a single card, which let us both recover quality and attack the MoE
plumbing:

\begin{itemize}
\tightlist
\item
  \textbf{Training-free recovery reproduces the down-anchor rule on a
  third model.} All-expert 2:4 on both projections is incoherent (PPL
  \textbf{111} vs stock \textbf{3.886}), but anchoring gate/up dense and
  sparsifying only the down projection recovers to \textbf{6.499 PPL --
  within 1.66x of stock, zero training}. The tax living in gate/up now
  holds across DeepSeek, GLM, \emph{and} gpt-oss.
\item
  \textbf{Fused monolith MoE.} A minimal-kernel pipeline (a
  counting-sort \texttt{moe\_route}; \texttt{gu\_swiglu} fusing GEMM1 +
  clamped-SwiGLU + MXFP4 quant in one launch; \texttt{seg\_bw} fusing
  down + bias + routing-weight into a coalesced epilogue; an
  inverse-index \texttt{moe\_combine} top-k) cuts the MoE to
  \textbf{\textasciitilde9.5 ms/layer, 2.2x} over the unfused path, all
  stages cos \textgreater= 0.996, and a same-invocation A/B measures
  \textbf{+6.0\% end-to-end serve}.
\end{itemize}

\textbf{Honest throughput.} The 2:4-sparse-FP4 GEMM beats vLLM's
\emph{real} \texttt{marlin\_gemm} by \textbf{2.15-3.04x} and the full
2-GEMM Marlin MoE by \textbf{1.81-2.24x} at gpt-oss/GLM/DeepSeek dims --
but Marlin is W4A16, a soft target. Against the real FP4 SOTA
(FlashInfer \texttt{cutlass} dense \texttt{mm\_fp4}, same invocation,
M=8192) our sparse kernel \textbf{ties on gpt-oss and loses 0.82-0.84x
on GLM/DeepSeek dims despite doing half the FLOPs}, consistent with the
dense verdict of Section 4. In serving, gpt-oss's uniquely \emph{small}
experts (2944x3072, \textasciitilde1016 tokens/expert) are Marlin's best
shape, and the per-layer W4A4 activation-quant Marlin skips makes stock
Marlin faster end-to-end. An exhaustive roofline settles that this is
silicon, not tuning: a register-only sparse-mma microbench
(\texttt{peak\_fed}) sustains \textbf{97.7\% of the sparse-mma peak}, so
the kernel is not feed-starved; the deficit is occupancy -- the
128-accumulator register file pins the high-reuse tile at 1 CTA/SM, and
SM120 has no \texttt{tcgen05}/TMEM to offload accumulators, while the
2:4-metadata and block-scale smem tax blocks the deeper pipeline dense
CUTLASS uses. A dense-throughput win on SM120 is physically unavailable.
gpt-oss's win is therefore the same as the other two models:
\textbf{only-deployed 2:4-sparse FP4, training-free capability recovery,
and \textasciitilde24\% weight-memory reduction at full sparsity} -- not
gpt-oss decode/prefill throughput, which lives on the large-expert
frontier models where the GEMM dominates the forward. Harnesses:
\texttt{harness/\{gptoss\_prep,fused\_moe,serve\_gptoss\}.py}.

\begin{center}\rule{0.5\linewidth}{0.5pt}\end{center}

\subsection{11. Related work}\label{related-work}

\textbf{FlashInfer \texttt{mm\_fp4}} is now the strongest dense FP4
baseline on SM120, and it moved while we worked. Its \texttt{auto}
selects \texttt{b12x} (a CUDA-13-only SM120/121 NVFP4 kernel) then
\texttt{cutlass}, and on the leaderboard (Section 4) these beat our
two-level dense by 1.35 to 2.2x. It is not a free win for the ecosystem:
the \texttt{cudnn} backend fails on every SM120 shape (cuDNN below
9.14), \texttt{b12x} collapses \textasciitilde2.2x at large serving
batch, \texttt{trtllm}/\texttt{cute-dsl} refuse SM120, and prebuilt
cubins historically shipped no sm\_120 targets (issue \#3294), so the
fast path needs CUDA-13 JIT. We benchmark against every one of its
backends rather than a single number.

\textbf{CUTLASS} ships dense (79b) and sparse (80b) NVFP4 GEMMs for
SM120 since 3.9.0. 80b is the only other sparse FP4 kernel in existence
and is our sparse baseline; the SM120 block-scaled path has documented
correctness and autotuner problems (issue \#3096: grouped GEMM garbage
output, TMA warp-specialized tactics failing to init), and, critically,
no library wraps \emph{any} sparse FP4 kernel in a deployment stack,
which is the gap the sparse path fills.

\textbf{Deployed inference on SM120.} The Marlin W4A16 fallback (dequant
FP4 to bf16 in-kernel, forfeiting the FP4 FLOPS) was reported around 50
tok/s on a 397B MoE, but for the dense Llama-3.1-8B NVFP4 checkpoint we
measure both vLLM 0.21 and SGLang 0.5 binding the \emph{native} NVFP4
W4A4 cutlass path on this card (Section 9: vLLM \texttt{modelopt\_fp4},
SGLang FlashInfer CUTLASS \texttt{fp4\_gemm}), not Marlin. Against those
real engines our through-kernel W4A4 accuracy matches (PPL 7.90 vs
7.97); the remaining gap is serving-stack engineering, not the kernel.

\textbf{Datacenter FP4 (B200, SM100, not our card).}
SGLang/FlashInfer/vLLM FP4 MoE reach roughly 1000 to 1260 TFLOPS using
\texttt{tcgen05}/UMMA that SM120 does not have. We do not compare our
SM120 numbers to these; the silicon differs.

\textbf{Quantization and sparsity recovery.} NVIDIA's NVFP4-QAD (arXiv
2601.20088) is the quant-recovery equivalent and recovers over 95\% of
FP accuracy on real Nemotron/Llama models; our sparse recovery is not
yet in that league. SparseGPT, and the adjacent SQ-format and SharQ
lines, are the sparse-plus-quant prior art we build on and must be
distinguished from; our specific move is retargeting the mask to
pair-granular 2:4 for the FP4 sparse path.

\begin{center}\rule{0.5\linewidth}{0.5pt}\end{center}

\subsection{12. Limitations}\label{limitations}

\begin{itemize}
\tightlist
\item
  \textbf{Serving-engine integration is built, graph-capturable, and
  beats production NVFP4 on decode; prefill still trails} (Section 9,
  Table C). quadbit's two-level sparse MLP runs inside vLLM (V1: paged
  attention, continuous batching, decode scheduler) as a
  \texttt{torch.library} custom op that vLLM's fullgraph compile +
  CUDA-graph capture include, NVFP4 for non-MLP, on a recovered
  Llama-3.1-8B-Instruct checkpoint (accuracy + speed on ONE checkpoint),
  with correct sparse accuracy under graphs (PPL 10.2709, not the 7.97
  dense value). A \textbf{split-K down projection}
  (\texttt{matmul\_sp\_sk}, 16→128 CTAs, down 109→56.5 µs) flips decode
  to a \textbf{win: +9.7/+7.2/+2.2\% at B=8/32/64} vs production NVFP4.
  \textbf{Prefill still trails \textasciitilde3-5\%} (−5.3 to −3.4\%);
  it uses the plain non-split-K down, which was never underfilled, and
  closing it needs a prefill-shape scheduling pass (stream-K /
  persistent tiling) that we have not yet built. The eager-vs-eager win
  (+5.6\% prefill, +23\% decode) was launch-overhead only and is a
  separate lever from the graph decode win. \textbf{Future work:} lift
  prefill to parity; graph-friendly SwiGLU/reduction fusion; accuracy
  recovery / hybrid sparsity.
\item
  \textbf{Sparse recovery loses to dense on accuracy} (Section 8). The
  deployable sparse-recovered 8B is 8.47 PPL through the two-level
  kernel, about 1.56 above dense W4A4 (6.91), and diverse-corpus data
  (C4) did not close it on the WikiText-2 test. Sparse is speed-only on
  accuracy grounds.
\item
  \textbf{The deploy gap is now closed.} The two-level sparse kernel
  (per-row and per-column fp32 global rescale in the epilogue) makes the
  deployed accuracy equal the trained fake-quant (8.47 == 8.47; the old
  single-level kernel would deploy the same checkpoint at 11.89). The
  fix costs 2 to 10\% of throughput and still beats CUTLASS 80b.
\item
  \textbf{Dense loses the SM120 FP4 race to FlashInfer.} The leaderboard
  (Section 4) shows FlashInfer \texttt{b12x}/\texttt{cutlass} beating
  our two-level dense by 1.35 to 2.2x. The dense kernel remains a useful
  zero-training W4A4 drop-in on the accuracy axis, but it is not the
  fastest dense FP4 on the platform. The win is sparse (the only
  deployed sparse FP4 GEMM, faster in wall-clock than even the best
  FlashInfer dense on prefill shapes), not dense.
\item
  \textbf{The Pareto win carries an accuracy tax.} The sparse speed
  advantage only pays off where a weight tolerates 2:4 pruning, and on
  8B that costs \textasciitilde1.56 PPL vs dense; the leaderboard result
  is a speed Pareto point, not a free lunch. End to end in serving the
  tax is a constant +2.3 PPL (10.27 vs 7.97 dense NVFP4), and it is the
  same across the whole request matrix.
\item
  \textbf{The sm\_120a block-scale mma is CUDA-12.8-only, forcing a
  staged build.} quadbit's \texttt{sm\_120a} block-scale mma
  (\texttt{kind::mxf4nvf4}/\texttt{block\_scale}/\texttt{scale\_vec::4X})
  is rejected by ptxas 13 and assembles only under CUDA \textless= 12.8,
  while FlashInfer's \texttt{b12x} needs CUDA 13, so the two cannot
  coexist in one container. The deployed path compiles the \texttt{.so}
  in a CUDA 12.8.1 image and ctypes-loads it into a CUDA 12.9 vLLM
  process, staged via a shared volume (Section 9,
  \texttt{docs/repro\_appendix.md}).
\item
  \textbf{int32 indexing caps a single GEMM.} The kernels pass
  \texttt{M}, \texttt{N}, \texttt{Klog} and compute all shared memory
  and global offsets as 32-bit \texttt{int}
  (\texttt{cuda/sparse\_fp4\_lib.cu}, \texttt{cuda/dense\_fp4\_lib.cu}),
  so a single launch is bounded to problems whose addressed element
  counts fit in int32. Real LLM linear shapes are far below this bound,
  but a 64-bit indexing pass would be required before the kernels could
  serve arbitrarily large fused shapes.
\item
  \textbf{Accuracy repair repairs PPL, not capability.} A four-way
  tournament (zero-runtime calibration, low-rank adapters, Wanda-pair
  mask repair, distillation) found only distillation moves the metric:
  it cuts the serving tax from 10.27 to 9.10 PPL while keeping every
  serving win, but the downstream-task accuracy (ARC-C/HellaSwag) stays
  \textasciitilde20 points below dense, so the capability loss from 2:4
  sparsity is not recovered (Section 9). Distillation buys perplexity,
  not capability; the honest open frontier is sparse capability
  recovery, not serving plumbing.
\item
  \textbf{No dense FP4 speed win.} quadbit dense loses the SM120 dense
  race to FlashInfer 1.35 to 2.2x and is retained only as a
  zero-training W4A4 accuracy drop-in, not a speed contribution.
\item
  \textbf{MoE decode occupancy not yet measured.} The segmented expert
  kernel is validated for correctness and scales at prefill routed-row
  counts (Section 10); at decode (few routed rows) its grid is small
  (the 16-CTA underfill regime of Section 8), and a split-K segmented
  variant -- the mechanical analogue of the single-MLP split-K down we
  already ship -- is future work, not yet built or measured. Because
  end-to-end DeepSeek-V4-Flash serving is ecosystem-blocked on SM120
  (Section 10), decode latency for this model could not be exercised
  regardless.
\item
  \textbf{Layerwise QAT does not recover the MoE gate/up tax;
  SparseGPT/distillation at MoE scale untried} (Section 10.2). We ran
  STE-QAT at MoE scale in two forms -- per-expert own-output, and global
  router-weighted combine fit to the dense routed aggregate -- and both
  fail to beat the one-shot sparse baseline (.7332 and .7259 vs .7363 on
  the 8-task battery), the tightest reconstruction giving the worst
  downstream; full-stack through-kernel QAT on dense Llama-3.1-8B is
  negative too (0.3967 \textless{} 0.4333). The gate/up 2:4 floor is a
  real capability property, avoided by structural placement, not
  repaired. What remains genuinely untried is calibrated
  SparseGPT/distillation on the experts (the dense-model levers of
  Section 8) at MoE scale, and true end-to-end QAT trainable
  \emph{through} the sparse kernel, which is infra-blocked (no
  differentiable DeepSeek/GLM load), not refuted.
\item
  \textbf{The MoE plugin path is graph-enabled and the dense-anchor
  bottleneck is removed} (Section 10, P4 + C1). P4 replaced the old
  host-syncing \texttt{torch.unique(...).tolist()} expert loop with a
  fixed-capacity device-routing path (\texttt{route\_fixed\_cap} /
  \texttt{\_route\_slot\_apply\_gs}), so the deployed sparse policies
  CUDA-graph-capture on SM120 (DeepSeek-D2, GLM route-slot D2, native
  SM120 DSA sparse-MLA, \texttt{drop=0}). C1 then removed the remaining
  decode-speed limit by delegating the dense anchored/grouped projection
  to FlashInfer's native grouped NVFP4 GEMM
  (\texttt{group\_gemm\_nvfp4\_nt\_groupwise}, opt-in
  \texttt{QB\_DENSE\_BACKEND=native\_nvfp4}) instead of the
  dequant-to-bf16 loop (\textbf{no custom dense grouped-GEMM was
  required}), so the captured DeepSeek-D2 path decodes faster than eager
  (5.82 vs 4.04 tok/s at matched PPL). This is a
  \textbf{same-model/same-policy speed/quality/memory Pareto result}
  against our own dequant-loop and eager paths, \textbf{not a
  production-wide decode-speed win} over other serving stacks. The
  remaining honest caveats are backend-scoped: the native delegate
  depends on FlashInfer availability and its swizzled NVFP4 scale
  layout, and the GLM graph rows are validated on a short held-out
  passage. Not a DSA, attention, memory, or loader blocker.
\item
  \textbf{The sparse MoE policy is not a decode-speed or decode-memory
  SOTA (Section 10.1, C2).} A direct same-harness board shows vLLM's
  native FlashInfer-CUTLASS fused NVFP4 MoE decodes 6-8x faster than the
  quadbit sparse route-slot D2 policy (DeepSeek 48.2 vs 6.0, GLM 33.8 vs
  5.4 tok/s captured) and uses less memory (D2's dual residency costs
  +26-27\% weight memory). The sparse-MoE value is quality-preserving
  structural sparsity + graph-enabled cross-architecture transfer + the
  Section 5 prefill/large-M kernel Pareto, \textbf{not} MoE decode
  speed; closing the decode gap needs a fused sparse grouped decode GEMM
  (new CUDA, identified but not built).
\item
  \textbf{GLM-5.2 requires 8x RTX PRO 6000.} At 432.9 GiB it does not
  fit on 2 or 4 cards; the smaller footprint DeepSeek-V4-Flash enjoyed
  (down-only at 2 GPUs) does not transfer. Route-slot dual residency
  fits on the 8-GPU host but drops KV capacity from 607k to 241k tokens
  (raw NVFP4 + 2:4 codes co-resident), so long-context KV pressure is
  the trade.
\item
  \textbf{GLM downstream is a smoke suite, not a full benchmark.} The
  GLM quality evidence is held-out PPL (all policies) plus a 4-task MC
  downstream comparison on the route-slot D2 policy only (Section 10).
  Full-size benchmarks, and downstream numbers for the down-only/gate-up
  GLM rows, are unmeasured; we make \textbf{no claim of exhaustive GLM
  downstream preservation}. The most complete downstream accounting of
  record remains DeepSeek's (full AVG across every policy).
\item
  \textbf{gpt-oss-20b is served and recovered, but does not win
  throughput} (Section 10.3). We serve gpt-oss end-to-end with the
  2:4-sparse experts live and recover quality training-free (111
  -\textgreater{} 6.499 PPL), and the fused monolith gives +6\% serve
  over our own unfused path; but stock Marlin W4A16 still serves gpt-oss
  faster, because its uniquely small experts (2944x3072,
  \textasciitilde1016 tok/expert) are Marlin's best shape and W4A4 pays
  a per-layer activation-quant Marlin skips. The GEMM math wins
  (2.15-3.04x vs real Marlin; tie/-18\% vs FlashInfer cutlass dense at
  M=8192), but on SM120 a dense-throughput serving win is physically
  unavailable (no \texttt{tcgen05}/TMEM; 1-CTA/SM occupancy on the
  high-reuse tile). gpt-oss's contribution is quality +
  \textasciitilde24\% memory + being the only deployed sparse-FP4 path,
  \textbf{not} its throughput; the throughput win lives on large-expert
  frontier models where the GEMM dominates the forward. The gpt-oss
  quality evidence is a held-out-passage PPL, not a downstream MC suite.
\end{itemize}

\begin{center}\rule{0.5\linewidth}{0.5pt}\end{center}

\subsection{13. Conclusion}\label{conclusion}

On SM120, FP4 is a real speedup but the ecosystem's coverage is uneven,
and it moved under us. On dense FP4 we no longer lead: FlashInfer's
CUDA-13 \texttt{b12x}/\texttt{cutlass} kernels beat our hand-written
two-level dense by 1.35 to 2.2x, and we report that plainly. What stands
is a Pareto point no shipping library provides. quadbit is the only
\emph{deployed} 2:4-sparse FP4 GEMM on SM120, its sparse kernel beats
CUTLASS 80b (the only other sparse kernel) on every shape, and in
wall-clock it beats even the best available \emph{dense} FP4 kernel
(FlashInfer \texttt{b12x}/\texttt{cutlass}) on every Llama-3-8B prefill
shape by 1.07 to 1.38x: if a weight can be 2:4-pruned, quadbit sparse is
the fastest way to run its FP4 GEMM on the platform. The accuracy axis
is honest too: the deployed dense W4A4 path costs +0.63 PPL with no
calibration, and sparse deploys at its trained accuracy through the
two-level kernel but stays \textasciitilde1.56 PPL behind dense, so the
sparse advantage is speed, conditioned on prunability. The contribution
is a hand-written kernel plus deployment stack that occupies a Pareto
corner (sparse FP4, deployed, fastest-for-prunable) that CUTLASS,
FlashInfer, SGLang, and vLLM leave empty.

The strongest results extend this beyond the dense single-GPU story into
cross-architecture sparse-policy transfer (Section 10). First, on a
large MoE the accuracy cost is not a fixed tax but a placement problem:
sparsifying only down projections in later layers, or only the
low-weight routed slots, preserves downstream capability
\textbf{training-free} on DeepSeek-V4-Flash (-0.29pt at 49\% of MoE
layers sparse; a per-expert weight repair, by contrast, fails). Second,
the same structural rules transfer to \textbf{GLM-5.2} served on 8x RTX
PRO 6000, whose Deep Sparse Attention runs natively on SM120: down-only
sparsity costs about half of gate/up there too, and the route-slot D2
policy costs +0.065 held-out PPL while preserving a 4-task downstream
smoke-suite average to within about one point of dense (.7508 vs .7603,
no task collapsing) -- so D2's small PPL cost is not masking a
downstream collapse. Third, the down-anchor rule reaches a \textbf{third
architecture, gpt-oss-20b}, which we serve end-to-end in vLLM with
training-free recovery (111 -\textgreater{} 6.499 PPL) and a fused
monolithic sparse MoE, the most complete serving of the three. The
deployed sparse policies now CUDA-graph-capture on SM120
(native-delegate captured DeepSeek-D2 decodes 1.44x its eager
reference), and we close the ``can training recover the gate/up tax?''
question in the negative across the whole layerwise family: per-expert
and global router-weighted MoE QAT both fail, the tightest
reconstruction giving the worst downstream, so the capability-preserving
move is structural placement, not weight repair. The honest ceilings
remain a full GLM downstream benchmark, a fused sparse grouped
\emph{decode} GEMM (the dense NVFP4 fused MoE is the SM120 decode SOTA),
and true end-to-end QAT through the sparse kernel (infra-blocked, not
refuted).

\begin{center}\rule{0.5\linewidth}{0.5pt}\end{center}

\subsection{Appendix: reproducibility}\label{appendix-reproducibility}

\begin{itemize}
\tightlist
\item
  \textbf{Kernels.} \texttt{cuda/matmul\_sp\_wide\_swz2.cu} (unit sparse
  2731k), \texttt{cuda/matmul\_sp\_full\_wide.cu} (deployable sparse
  2116k), \texttt{cuda/matmul\_fp4\_pp\_bf16.cu} (dense 1503k),
  \texttt{cuda/dense\_fp4\_lib.cu}, \texttt{cuda/sparse\_fp4\_lib.cu}
  (PyTorch-callable plus fused quantizer).
\item
  \textbf{Probes.} \texttt{mma\_peak}, \texttt{tma\_bw}, \texttt{smemq},
  \texttt{sp\_*\_probe}, \texttt{verify\_*}, \texttt{pack\_verify},
  \texttt{pack\_accuracy}.
\item
  \textbf{Harnesses.} \texttt{bench\_vs\_bf16.py} (throughput),
  \texttt{cutlass\_fp4.py} / \texttt{cutlass\_sparse.py} /
  \texttt{cutlass\_shapes.py} (CUTLASS baselines),
  \texttt{quadbit\_linear.py} (drop-in), \texttt{recovery\_worth.py}
  (dense W4A4 accuracy), \texttt{sparsegpt\_pair.py} (one-shot),
  \texttt{finetune\_pair.py} (recovery), \texttt{sensitivity\_sparse.py}
  (hybrid sparse-placement sweep, training-free negative result),
  \texttt{vllm\_nvfp4.py} (vLLM SM120 NVFP4 smoke test).
\item
  \textbf{Full chronological build log.} Memory
  \texttt{quadbit-raw-ptx.md}.
\item
  \textbf{Exact commands, model ids, checkpoints, and the staged .so
  build recipe.} \texttt{docs/repro\_appendix.md}.
\item
  \textbf{Figure plan (what each figure shows and its data source).}
  \texttt{docs/figure\_plan.md}.
\item
  \textbf{Claims checklist (every claim to its evidence and status).}
  \texttt{docs/claims\_checklist.md}.
\end{itemize}

\end{document}
